\documentclass[11pt,letterpaper,oneside]{article}
\usepackage[T1]{fontenc}
\usepackage[utf8]{inputenc}
\usepackage{lmodern}
\usepackage[margin=1in,headheight=15pt]{geometry}
\usepackage{amsmath,amssymb,amsthm,mathtools,mathrsfs}
\usepackage{microtype,booktabs,longtable,array,enumitem}
\usepackage[dvipsnames]{xcolor}
\usepackage{fancyhdr}
\usepackage[colorlinks=true,linkcolor=MidnightBlue,citecolor=MidnightBlue,urlcolor=MidnightBlue]{hyperref}
\usepackage[nameinlink,noabbrev]{cleveref}
\hypersetup{pdftitle={Surcomplex Polynomial Algebra: Ordered Geometry, Valuative Root Trees, and Residue Duality},pdfauthor={Research exposition prepared with ChatGPT}}
\pagestyle{fancy}\fancyhf{}
\fancyhead[L]{\small Surcomplex polynomial algebra}
\fancyhead[R]{\small Geometry, valuations, and duality}
\fancyfoot[C]{\thepage}
\renewcommand{\headrulewidth}{0.3pt}
\setlength{\parskip}{0.3em}\setlength{\parindent}{1.25em}
\setlist[enumerate]{leftmargin=*,itemsep=0.2em,topsep=0.35em}
\numberwithin{equation}{section}
\newtheorem{theorem}{Theorem}[section]
\newtheorem{lemma}[theorem]{Lemma}
\newtheorem{proposition}[theorem]{Proposition}
\newtheorem{corollary}[theorem]{Corollary}
\theoremstyle{definition}
\newtheorem{definition}[theorem]{Definition}
\newtheorem{example}[theorem]{Example}
\theoremstyle{remark}\newtheorem{remark}[theorem]{Remark}
\newcommand{\No}{\mathbf{No}}\newcommand{\SC}{\mathbf{SC}}
\newcommand{\C}{\mathbb C}\newcommand{\R}{\mathbb R}\newcommand{\Q}{\mathbb Q}
\newcommand{\N}{\mathbb N}\newcommand{\Z}{\mathbb Z}
\newcommand{\OO}{\mathcal O}\newcommand{\mm}{\mathfrak m}
\newcommand{\HH}{\mathcal H_\Gamma}\newcommand{\Hp}{\mathcal H^+_\Gamma}
\newcommand{\sh}{^{\sharp}}
\newcommand{\supp}{\operatorname{supp}}\newcommand{\st}{\operatorname{st}}
\newcommand{\lc}{\operatorname{lc}}\newcommand{\ord}{\operatorname{ord}}
\newcommand{\NF}{\operatorname{NF}}\newcommand{\Tr}{\operatorname{Tr}}
\newcommand{\Res}{\operatorname{Res}}\newcommand{\Disc}{\operatorname{Disc}}
\newcommand{\rank}{\operatorname{rank}}\newcommand{\sig}{\operatorname{sig}}
\newcommand{\sgn}{\operatorname{sgn}}\newcommand{\Hom}{\operatorname{Hom}}
\newcommand{\conv}{\operatorname{conv}}\newcommand{\Bez}{\operatorname{Bez}}
\newcommand{\id}{\operatorname{id}}\newcommand{\re}{\operatorname{Re}}
\newcommand{\im}{\operatorname{Im}}\newcommand{\hoint}{\mathop{\oint}\nolimits^{\mathrm H}}
\newcommand{\Bcl}{B_{\geq}}\newcommand{\Bop}{B_{>}}
\newcommand{\abs}[1]{\left|#1\right|}
\newcolumntype{L}[1]{>{\raggedright\arraybackslash}p{#1}}
\begin{document}
\hypersetup{pageanchor=false}
\begin{titlepage}
\centering
\vspace*{0.65in}
{\large\scshape A theorem-driven continuation\par}
\vspace{0.45in}
{\Huge\bfseries Surcomplex\par Polynomial Algebra\par}
\vspace{0.3in}
{\LARGE Ordered Geometry,\par Valuative Root Trees,\par and Residue Duality\par}
\vspace{0.4in}
{\large September 21, 2026\par}
\vspace{0.35in}
\begin{minipage}{0.94\textwidth}
\begin{abstract}
We develop finite polynomial algebra over the surcomplex class
$\SC=\No[i]$, using algebraically closed, set-sized Hahn fields to make every
factorization and quotient an ordinary set-sized construction. The theory has
three interacting structures: algebraic closedness, ordered complex geometry,
and an arbitrary-rank valuation with residue field $\C$.

We prove surcomplex forms of Gauss--Lucas, polynomial Rouch\'e, Hermite's
signature criterion, Newton-polygon root counting, coprime Hensel lifting,
and exact Rolle counts in valuative disks. A residual derivative polynomial
locates the critical points between the children of each root cluster.
Discriminants become weighted sums on the cluster tree, and an explicit
coefficient-precision theorem preserves all pairwise root separations while
giving the leading displacement of each root.

For polynomial systems with coordinate-power highest terms, we construct a
global finite-free quotient over any commutative coefficient ring, including
Hahn-coherent parameter rings. A total-degree argument replaces the
positive-support hypothesis needed in the supplied analytic manuscripts.
An explicit multivariate B\'ezoutian is the duality kernel; on its diagonal it
is the Jacobian, giving an exact trace--residue formula through collisions.
For positive lower-degree perturbations this also identifies the polynomial
restriction of the contour residue in the supplied research manuscript.
All major assertions are proved, with examples and executable exact checks.
\end{abstract}
\end{minipage}
\vfill
\begin{minipage}{0.94\textwidth}\small
\textbf{Provenance.} The two supplied manuscripts are treated as unpublished
source material. Classical algebra, real-closed-field transfer, valued-field
root theory, and B\'ezoutian duality have substantial prior literature. This
article is a proved synthesis and extension of the supplied framework, not a
claim to solve a named open conjecture or establish priority. The proofs have
not been independently refereed or formally verified.
\end{minipage}
\end{titlepage}
\hypersetup{pageanchor=true}
\pagenumbering{roman}
\setcounter{tocdepth}{1}\tableofcontents
\clearpage
\pagenumbering{arabic}

\section{Scope, source relationship, and principal results}\label{sec:scope}
A surcomplex number is $x+iy$ with $x,y\in\No$ and $i^2=-1$. The
algebraic closedness of this class field is established background, not a
new fundamental theorem of algebra \cite{Conway,Gonshor,vdDE,Poonen}.
The real question is which \emph{additional structures} a polynomial theorem
uses. An identity of resultants uses only a field. Convexity uses conjugation
and the order on its real part. A Newton polygon uses a valuation. A contour
argument uses an analytic category that is not supplied by algebraic
closedness alone.

The supplied \emph{Surcomplex Analysis} manuscript \cite{SourceAnalysis}
distinguishes fine-local analytic germs, Hahn-coherent functions on ordinary
complex domains thickened by infinitesimals, and all-scale coherence. Its
preparation and residue theorems already handle infinitesimally displaced
zeros. The supplied \emph{Hartogs Coherence, Finite Maps, and Residue Duality}
manuscript \cite{SourceResearch} develops parameter division, coupled analytic
systems, and residue pairings. We retain their notation $t^\gamma=
\omega^{-\gamma}$, their use of set-sized Hahn workspaces, and their meaning
of strong summability.

The present article changes the central object from a coherent analytic
function to a \emph{finite polynomial}. This permits evaluations at infinite
as well as finite surcomplex points, finite elimination procedures, and
arbitrary coefficient sizes. It also requires care: the global polynomial
quotient need not be the finite-monad analytic quotient.

\begin{example}[A global root invisible in a finite monad]\label{ex:escape}
Let $t=\omega^{-1}$ and $F(z)=z+tz^2=z(1+tz)$. As a global polynomial it
has the two roots $0$ and $-t^{-1}$, and $K[z]/(F)$ has dimension two.
On the thickening of any fixed ordinary disk about zero, $1+tz$ is a
Hahn-coherent unit, with strong inverse $\sum_{k\geq0}(-tz)^k$.
The corresponding analytic quotient has rank one and sees only the root
zero. Neither statement contradicts the other. In particular, positive
Hahn support of a perturbation does \emph{not} by itself preserve the global
polynomial degree or the total number of roots over all of $\SC$.
\end{example}

\subsection{The main theorem packages}
The central valuation results hold in every divisible Hahn workspace.
Its exponent group need not be Archimedean or have rank one.
\begin{center}\small
\begin{tabular}{@{}L{0.27\textwidth}L{0.61\textwidth}@{}}
\toprule
Package & Principal conclusion\\
\midrule
Ordered root geometry & Critical points are surreal convex combinations of
roots; real-root counts are signatures of finite trace matrices.\\[0.45em]
Weighted root counting & The degree and zero multiplicity of an initial
polynomial count roots in closed and open valuative disks.\\[0.45em]
Critical-point clusters & A disk containing $m\geq1$ roots contains exactly
$m-1$ roots of the derivative, with multiplicity.\\[0.45em]
Precision and stability & For integral simple roots $r_i$, coefficient
precision greater than $\max_i(vP'(r_i)+\delta_i)$ preserves the entire
root tree, where $\delta_i$ is the nearest-neighbor valuation.\\[0.45em]
Polynomial complete intersections & Systems $F_i=x_i^{d_i}+E_i$ with
$\deg E_i<d_i$ have a rectangular basis of size $\prod_i d_i$;
B\'ezoutian duality gives $\Tr(m_H)=\lambda_F(J_FH)$.\\
\bottomrule
\end{tabular}
\end{center}

We provide a direct proof of the last package over an arbitrary commutative
ring. This covers non-Noetherian integral Hahn coefficient rings and is
particularly useful for comparison with \cite{SourceResearch}. The degree
hypothesis is a genuine restriction, not a disguised assertion about all
isolated complete intersections.

\subsection{What is and is not claimed}
Most underlying principles are classical. Our contribution here is the
precise surcomplex formulation, simultaneous control of order and valuation,
explicit root-tree precision, and a direct polynomial comparison with the
supplied analytic residues. The cited literature includes non-Archimedean
root location \cite{ChoiLee}, real-algebraic complex root counting
\cite{Eisermann,Swan}, and multivariate B\'ezoutian duality
\cite{BCRS,BMP}. No theorem is declared new merely because ``surcomplex'' is
placed in its title.

No assertion below treats a proper-class-indexed sum, an infinite-degree
object as a polynomial, or a strong Hahn sum as the limit of its partial
sums in the full fine topology. The latter distinction is essential in both
supplied manuscripts and remains essential here.

\section{Foundations: fields, order, valuation, and supports}\label{sec:foundations}
\subsection{Set-sized workspaces}
Work in a class set theory such as NBG with global choice. All individual
polynomials have ordinary finite degree. If $\Gamma\subseteq\No$ is a
set-sized divisible additive subgroup, write
\[
 R_\Gamma=\R((t^\Gamma)),\qquad
 K_\Gamma=\C((t^\Gamma))=R_\Gamma[i],\qquad
 t^\gamma=\omega^{-\gamma}.
\]
A series $A=\sum_{\gamma\in S}a_\gamma t^\gamma$ has a well-ordered
set $S\subseteq\Gamma$ as support. Conway normal forms embed these fields
in $\No$ and $\SC$. A finite, or any set-sized, collection of coefficients
lies in such a workspace: take the divisible subgroup generated by all
exponents in their supports.

We use the classical algebraic-closure theorem for Hahn fields with divisible
value group and algebraically closed residue field; see
\cite[Corollary 4]{Poonen}. Hence $K_\Gamma$ is algebraically closed.
The ordered field $R_\Gamma$ is real closed. For instance, its complexification
is algebraically closed and it is formally real, so the standard real-closed
field criterion applies. All subsequent root factorizations are performed
in a field $K=K_\Gamma$, not by an unexplained proper-class operation.

\begin{proposition}[Workspace independence]\label{prop:workspace}
Let $P\in K[z]$ be nonzero. Every surcomplex root of $P$ belongs to $K$.
More generally, if $A=K[x_1,\ldots,x_s]/I$ is finite-dimensional, every
surcomplex point annihilating $I$ has all coordinates in $K$. Its local
multiplicity is unchanged by enlargement to another algebraically closed
Hahn workspace.
\end{proposition}
\begin{proof}
Factor $P=a\prod_j(z-r_j)$ in $K[z]$. At a zero in any larger field, one
linear factor must vanish. For a finite algebra, multiplication by $x_i$
satisfies its characteristic polynomial. At a point annihilating $I$ its
$i$th coordinate is a root of that polynomial and therefore lies in $K$.
The finite algebra is a product of local Artinian algebras with residue
field $K$. Scalar extension preserves each factor's dimension; its extended
maximal ideal is still nilpotent and its residue field is the new field,
so the factor remains local. This proves the multiplicity assertion.
\end{proof}

\subsection{Modulus and the natural valuation}
For $z=x+iy$ set
\[
 \overline z=x-iy,\qquad |z|=\sqrt{x^2+y^2}.
\]
The algebraic proofs of $|zw|=|z||w|$ and $|z+w|\leq|z|+|w|$ work over
any real closed field. For a nonzero Hahn series define
\[
 v(A)=\min\supp(A),\qquad \lc(A)=a_{v(A)},\qquad v(0)=+\infty.
\]
Then
\[
 v(AB)=v(A)+v(B),\quad
 v(A+B)\geq\min(v(A),v(B)),\quad v(|A|)=v(A).
\]
Every nonzero ordinary complex number has valuation zero. In particular,
\begin{equation}\label{eq:integersunit}
 v(m)=0\qquad(m\in\Z\setminus\{0\}).
\end{equation}
This residue-characteristic-zero fact is responsible for the exact
critical-point counts later in the article.

The valuation ring and maximal ideal are
\[
 \OO_K=\{z:v(z)\geq0\},\qquad \mm_K=\{z:v(z)>0\}.
\]
They consist respectively of finite elements and infinitesimals.
Reduction is the standard-part map
$\st:\OO_K\to\C$, taking the coefficient at exponent zero. Thus every
finite $z$ is uniquely $c+\varepsilon$, with $c\in\C$ and
$\varepsilon\in\mm_K$.

Two different geometries must be kept separate. An \emph{ordered disk} is
$\{z:|z-a|\leq R\}$ for a positive surreal $R$. The \emph{valuative disks}
used below are
\begin{equation}\label{eq:balls}
 \Bcl(a,\gamma)=\{z:v(z-a)\geq\gamma\},\qquad
 \Bop(a,\gamma)=\{z:v(z-a)>\gamma\}.
\end{equation}
For example, $\Bcl(0,0)$ consists of \emph{all} finite elements, not just
the ordered unit disk. The scale $t^\gamma$ specifies an Archimedean size
class, not a bound on its ordinary leading coefficient.

\subsection{Strong sums and finite polynomial degree}
A family of Hahn series is strongly summable if the union of its supports
is well ordered and every exponent occurs in only finitely many members.
The sum is then defined by finite coefficientwise addition. We use the
Hahn--Neumann lemma in the form recorded in both supplied manuscripts:
if $S,T$ are well ordered, then $S+T$ is well ordered with finite
representation multiplicities; if $E\subseteq\Gamma_{>0}$ is well ordered,
then its finite-word monoid $E^*$ is well ordered and each exponent has
only finitely many word representations, up to the finite permutations
of each word \cite{Poonen,SourceAnalysis}.

\begin{remark}[The order of the two constructions matters]\label{rem:degree}
$K[z]$ consists of expressions of bounded finite degree in $z$, although
each coefficient can have an infinite Hahn support. In contrast,
$\C[z]((t^\Gamma))$ permits unbounded polynomial degrees among its Hahn
coefficients. The formal expression $\sum_{k\geq0}t^kz^k$ belongs to the
latter ring when $1\in\Gamma$, but its substitution at $z=t^{-1}$ is the
nonsummable family $1+1+\cdots$. It is not a global surcomplex polynomial.
All global polynomial evaluations in this article are finite sums.
\end{remark}

\section{Finite algebra: factorization, elimination, and interpolation}\label{sec:algebra}
\subsection{The fundamental theorem and Euclidean algebra}
\begin{theorem}[Finite polynomial algebra over $\SC$]\label{thm:fta}
Every nonzero degree-$n$ polynomial with surcomplex coefficients has exactly $n$
roots counted with multiplicity, all in a single algebraically closed Hahn
workspace. Euclidean division, monic greatest common divisors, B\'ezout
identities, unique factorization, and the Chinese remainder theorem hold
for any specified finite collection of such polynomials.
\end{theorem}
\begin{proof}
Choose $K$ containing all coefficients. Algebraic closedness produces one
root; division by its monic linear factor reduces the degree by one.
Induction gives a complete factorization, unique up to ordering.
The ordinary leading-term subtraction algorithm divides one polynomial by
another in finitely many steps. Its successive nonzero remainder degrees
strictly decrease, proving termination of the Euclidean algorithm and
its B\'ezout identity. Coprime-factor Chinese remaindering follows by using
that identity to construct complementary idempotents. All operations
remain in $K$, and \cref{prop:workspace} excludes extra roots outside it.
\end{proof}

For $P=a_nz^n+\cdots+a_0=a_n\prod_{j=1}^n(z-r_j)$, Vieta's formulas
and the Newton identities are therefore unchanged. For example, after
making $P$ monic and writing its coefficients as $a_{n-1},\ldots,a_0$,
the sums $s_k=\sum_jr_j^k$ satisfy
\[
 s_k+a_{n-1}s_{k-1}+\cdots+a_{n-k+1}s_1+k a_{n-k}=0
 \quad(1\leq k\leq n).
\]
They follow by expanding the product, or by comparing coefficients in
$P'/P=\sum_j(z-r_j)^{-1}$ at infinity. This is a formal Laurent identity;
no limiting process at an ``infinite point'' is required.

The gcd $\gcd(P,P')$ records repeated roots. In characteristic zero,
$P/\gcd(P,P')$ is the monic squarefree polynomial with the same distinct
zeros, after normalization. Repeated gcd operations also give the usual
squarefree decomposition, separating each multiplicity.

\subsection{Resultants, discriminants, and multiplication matrices}
Let $P$ be monic of degree $n$, and let $C_P$ be the matrix of multiplication
by $z$ on $A_P=K[z]/(P)$ in the basis $1,z,\ldots,z^{n-1}$.
For every polynomial $Q$,
\begin{equation}\label{eq:resultantnorm}
 \det Q(C_P)=\prod_{j=1}^n Q(r_j)=\Res(P,Q).
\end{equation}
Indeed, the local Chinese remainder factors have the form
$K[u]/(u^{m_r})$, with $z=r+u$. Multiplication by $Q$ on such a factor
is triangular with diagonal $Q(r)$, repeated $m_r$ times. This proves
\eqref{eq:resultantnorm} and its version for a nonmonic $P$, whose
resultant is $a_n^{\deg Q}\prod_jQ(r_j)$.

The determinant of the Sylvester matrix gives the same resultant. One way
to see its zero criterion directly is to view that matrix as the map
$(U,V)\mapsto UP+VQ$ with the usual degree bounds. It is singular exactly
when $P$ and $Q$ share a nonconstant factor. Multiplicativity and the
linear-factor calculation then yield the displayed normalization.
Consequently the familiar elimination test remains exact even when the
resultant is infinitesimal or infinite.

For $n\geq2$ the discriminant is
\begin{equation}\label{eq:disc}
 \Disc(P)=a_n^{2n-2}\prod_{i<j}(r_i-r_j)^2
 =(-1)^{n(n-1)/2}a_n^{-1}\Res(P,P').
\end{equation}
It vanishes precisely at a repeated root. A tiny nonzero discriminant
is not zero: it measures close, but still algebraically distinct, roots.

\subsection{Hermite interpolation and rational partial fractions}
Suppose $r_1,\ldots,r_s$ are distinct and $m_1+\cdots+m_s=n$. Then
\begin{equation}\label{eq:hermiteCRT}
 K[z]\Big/\prod_{j=1}^s(z-r_j)^{m_j}
 \ \cong\ \prod_{j=1}^s K[u]/(u^{m_j}),
 \qquad H\longmapsto H(r_j+u)\bmod u^{m_j}.
\end{equation}
The ideals are pairwise coprime, so this is the Chinese remainder theorem.
It states that arbitrary prescribed jets
$H^{(k)}(r_j)$, $0\leq k<m_j$, determine a unique polynomial of degree
less than $n$. The factorials needed to convert Taylor coefficients to
derivatives are invertible.

When all $m_j=1$, the interpolation basis is
\[
 L_j(z)=\frac{P(z)}{(z-r_j)P'(r_j)},\qquad L_j(r_i)=\delta_{ij}.
\]
The denominators may be arbitrarily small, so the interpolation coefficients
may be arbitrarily large. The formula is nonetheless exact. Division and
\eqref{eq:hermiteCRT} similarly give every rational function a unique
polynomial part and finitely many terms $c_{j,k}(z-r_j)^{-k}$. This is the
algebraic foundation for the residue theory in \cref{sec:uniresidue}.

\section{Ordered geometry of roots and critical points}\label{sec:geometry}
All convex combinations in this section use nonnegative \emph{surreal}
weights, not just ordinary real weights. For a finite set $S\subseteq K$,
\[
 \conv_{R_\Gamma}(S)=
 \left\{\sum_{r\in S}\theta_r r:
    \theta_r\in R_\Gamma,\ \theta_r\geq0,\ \sum_r\theta_r=1\right\}.
\]
This is a finite-dimensional algebraic notion of convexity. No topological
compactness of a proper-class set is involved.

\begin{theorem}[Surcomplex Gauss--Lucas with explicit weights]\label{thm:gausslucas}
Let $P$ be nonconstant and let $w$ be a zero of $P'$. Then $w$ lies in the
surreal convex hull of the roots of $P$. If $P(w)\ne0$ and the roots are
listed with multiplicity, then
\begin{equation}\label{eq:barycentric}
 w=\sum_{j=1}^n\theta_j r_j,
 \qquad
 \theta_j=\frac{|w-r_j|^{-2}}{\sum_{k=1}^n|w-r_k|^{-2}}>0.
\end{equation}
\end{theorem}
\begin{proof}
A critical point that is also a root already belongs to the hull. Otherwise
\[
 0=\frac{P'(w)}{P(w)}=\sum_j\frac1{w-r_j}
 =\sum_j\frac{\overline w-\overline r_j}{|w-r_j|^2}.
\]
Conjugate this identity and divide by the strictly positive sum of the
real denominators' inverses. This gives \eqref{eq:barycentric}.
\end{proof}

Every ordered disk and every ordered half-plane is convex by the triangle
inequality and linearity of real part. Thus a disk or half-plane containing
all the roots contains all the critical points. Iteration gives the same
statement for every nonzero higher derivative. For a real-rooted polynomial,
every derivative is real-rooted.

\begin{proposition}[Weighted incomplete polynomials]\label{prop:weighted}
For $r_1,\ldots,r_n\in K$ and weights $\lambda_j\geq0$ in $R_\Gamma$
with $\sum_j\lambda_j>0$, every root of
\[
 Q_\lambda(z)=\sum_{j=1}^n\lambda_j
       \prod_{k\ne j}(z-r_k)
\]
belongs to $\conv_{R_\Gamma}\{r_1,\ldots,r_n\}$. Conversely, as
$\lambda$ varies, the union of these root sets is exactly that convex hull
when $n\geq2$.
\end{proposition}
\begin{proof}
At a root distinct from all $r_j$, divide by $\prod_j(z-r_j)$ and repeat
the preceding proof with weights $\lambda_j/|z-r_j|^2$.
For the converse let $w=\sum_j\theta_jr_j$, with
$\theta_j\geq0$ and $\sum_j\theta_j=1$. If $w$ is not a listed root,
choose $\lambda_j=\theta_j|w-r_j|^2$. Then
\[
 \sum_j\frac{\lambda_j}{w-r_j}
 =\sum_j\theta_j(\overline w-\overline r_j)=0.
\]
At a listed root $r_i$, choose a nonzero weight supported at an index
$j\ne i$. Its incomplete product contains the factor $z-r_i$.
\end{proof}
This is an ordered-field form of a classical extension of Gauss--Lucas;
compare the complex result in \cite{Zhang}. The proof exhibits precisely
why it survives for surreal coefficients.

\begin{proposition}[A Cauchy bound at arbitrary surreal size]\label{prop:cauchy}
Every root of a monic polynomial of positive degree
$P=z^n+a_{n-1}z^{n-1}+\cdots+a_0$ satisfies
\[
 |z|\leq1+\max_{0\leq j<n}|a_j|.
\]
\end{proposition}
\begin{proof}
Put $M=\max_j|a_j|$. If $r=|z|>1+M$, then
\[
 \left|\sum_{j<n}a_jz^j\right|
 \leq M\sum_{j<n}r^j
 =M\frac{r^n-1}{r-1}<r^n.
\]
The leading term cannot be cancelled. The geometric sum is finite, and the
argument is valid for an infinite surreal $r$ as well as a finite one.
\end{proof}

\section{Real-algebraic counting and an ordered-disk Rouch\'e theorem}\label{sec:transfer}
\subsection{The exact scope of real-closed-field transfer}
We use Tarski's classical principle: a first-order sentence in the language
of ordered fields is true in one real closed field if and only if it is
true in every real closed field \cite{Swan}. Complex variables can be
replaced by pairs of real variables, conjugation is polynomial in these
coordinates, and inequalities involving a modulus can be squared when
both sides are nonnegative.

This is a theorem about each \emph{fixed finite degree}. It does not
quantify over an infinite sequence of polynomial coefficients, the
standard natural numbers as a definable set, arbitrary subsets, or all
fine-continuous functions. For finite data in $\SC$, apply transfer in
$R_\Gamma$ and then use \cref{prop:workspace}. This avoids treating a
proper class as an ordinary first-order structure.

\begin{theorem}[Polynomial Rouch\'e on an ordered surreal disk]\label{thm:orderedrouche}
Let $P,Q\in\SC[z]$, let $a\in\SC$, and let $R\in\No_{>0}$. Suppose
\[
 |Q(z)-P(z)|<|P(z)|\quad\text{whenever }|z-a|=R.
\]
Then $P$ and $Q$ have the same number of roots in $|z-a|<R$, counted
with multiplicity, and neither has a root on the boundary.
\end{theorem}
\begin{proof}
Choose a real closed workspace containing the real and imaginary parts
of all coefficients, $a$, and $R$. Fix the two actual polynomial degrees.
For each possible count $0\leq k\leq\max(\deg P,\deg Q)$, saying
that a polynomial has exactly $k$ roots in the disk is a first-order
statement: quantify a complete factorization into finitely many linear
factors and require exactly $k$ of the listed factors' roots to satisfy
the strict disk inequality. Repetitions in that factorization count
multiplicity. The boundary inequality is also first order, using squared
moduli and a universal quantifier over two real coordinates.

For these fixed degrees the classical complex polynomial Rouch\'e theorem
asserts the resulting implication over $\R$. Tarski's principle gives
it over the workspace. Boundary nonvanishing follows at once from the
strict inequality, and workspace independence gives the conclusion in
$\SC$. No real-parameter fine-continuous circle or contour integral has
been introduced.
\end{proof}

This is a transfer theorem, not a new analytic proof. Real-algebraic
approaches via Sturm chains and algebraic winding numbers give a more
algorithmic treatment of complex polynomial counting over real closed
fields \cite{Eisermann}. The valuative Rouch\'e theorem proved below has
different hypotheses and admits a direct initial-polynomial proof.

\subsection{Hermite trace forms count real roots}
Let $R$ be a real closed field, $P\in R[z]$ monic of degree $n$, and
$A=R[z]/(P)$. For $q\in A$, define the symmetric bilinear form
\begin{equation}\label{eq:hermiteform}
 T_q(f,g)=\Tr_A(m_{qfg}).
\end{equation}
A symmetric matrix over a real closed field is diagonalizable by
congruence. Its signature is the number of positive diagonal entries
minus the number of negative ones, with zero entries ignored. Sylvester's
law of inertia makes this independent of the diagonalization.

\begin{theorem}[Hermite's signature criterion over surreal real fields]\label{thm:hermite}
The signature of \eqref{eq:hermiteform} is
\begin{equation}\label{eq:hermitesignature}
 \sig T_q=\sum_{\substack{r\in R\\P(r)=0}}\sgn q(r),
\end{equation}
where the sum runs over \emph{distinct} real roots. In particular,
$\sig T_1$ is their number. The rank of $T_1$ is the number of distinct
roots over $R[i]$, and $T_1$ is positive semidefinite if and only if every
root is real.
\end{theorem}
\begin{proof}
Factor $P$ over $R$ into powers of distinct real linear factors and
irreducible quadratic factors. Chinese remaindering makes the summands
orthogonal. On a local factor at a real root of multiplicity $m$, the
trace of multiplication by $h$ is $m h(r)$: its nilpotent part is
triangular with zero diagonal. Hence the form factors through the
one-dimensional residue field and is $m q(r)uv$ there. Its signature is
$\sgn q(r)$.

For a nonreal conjugate pair, the residue field is $R[i]$. If $q(r)=a+ib$,
its residue form, up to the positive factor $m$, has matrix
\[
 \begin{pmatrix}2a&-2b\\-2b&-2a\end{pmatrix}.
\]
If $q(r)\ne0$, its determinant is $-4(a^2+b^2)<0$, so its signature is
zero and its rank is two. If $q(r)=0$, it is zero. Nilpotent directions
lie in the kernel in both cases. Adding the signatures proves
\eqref{eq:hermitesignature}; taking $q=1$ gives the other assertions.
\end{proof}

This classical signature method is related to the residue and B\'ezoutian
forms discussed in \cite{BCRS}; it is not the nondegenerate residue pairing
of \cref{sec:uniresidue}. It gives a practical finite-matrix criterion:
its entries in the monomial basis are $\Tr(q(C_P)C_P^{i+j})$.
For $a<b$ with $P(a)P(b)\ne0$, put $q(z)=(z-a)(b-z)$. Then
\begin{equation}\label{eq:intervalcount}
 \#\{r\in(a,b):P(r)=0\}
 =\frac{\sig T_1+\sig T_q}{2},
\end{equation}
again counting distinct roots. Apply this separately to the squarefree
multiplicity factors to obtain a count with multiplicity.

\section{Initial polynomials and exact root counts at every scale}\label{sec:newton}
Fix an algebraically closed Hahn workspace $K=K_\Gamma$. Write $v(0)=+\infty$.
Given a nonzero polynomial $P$, a center $a\in K$, and $\gamma\in\Gamma$,
expand
\[
 P(a+Z)=\sum_{j=0}^n b_j Z^j.
\]
Define its weighted Gauss value and initial polynomial by
\begin{align}
 g_P(a,\gamma)&=\min_{b_j\ne0}\{v(b_j)+j\gamma\},\label{eq:gaussvalue}\\
 I_{a,\gamma}(P)(X)&=\st\!\left(t^{-g_P(a,\gamma)}P(a+t^\gamma X)\right)
 \in\C[X].\label{eq:initial}
\end{align}
The standard part in \eqref{eq:initial} is taken coefficientwise. All
coefficients inside the parentheses have nonnegative valuation, and at
least one is a unit, so $I_{a,\gamma}(P)\ne0$. This definition uses no
limiting process and is meaningful even for nonreal surreal exponents.

\begin{lemma}[Multiplicativity of initial polynomials]\label{lem:gaussproduct}
For nonzero $P,Q\in K[z]$,
\[
 g_{PQ}(a,\gamma)=g_P(a,\gamma)+g_Q(a,\gamma),\qquad
 I_{a,\gamma}(PQ)=I_{a,\gamma}(P)I_{a,\gamma}(Q).
\]
\end{lemma}
\begin{proof}
Normalize each polynomial as in \eqref{eq:initial}. Their product has
integral coefficients and has nonzero reduction, because $\C[X]$ is an
integral domain. Its least coefficient valuation is therefore zero.
Undoing the two normalizing monomials proves both identities.
\end{proof}

\begin{theorem}[Root counting by one initial polynomial]\label{thm:initialcount}
Let $P(z)=a_n\prod_{i=1}^n(z-r_i)$, with roots repeated according to
multiplicity. Put $\lambda_i=v(r_i-a)$, allowing $\lambda_i=+\infty$.
Then
\begin{equation}\label{eq:gaussroots}
 g_P(a,\gamma)=v(a_n)+\sum_{i=1}^n\min(\gamma,\lambda_i).
\end{equation}
For a nonzero complex scalar $c_*$,
\begin{equation}\label{eq:initialroots}
 I_{a,\gamma}(P)(X)
 =c_*X^{N_>}\prod_{\lambda_i=\gamma}
 \left(X-\st\frac{r_i-a}{t^\gamma}\right),
 \qquad N_>=\#\{i:\lambda_i>\gamma\}.
\end{equation}
Consequently, writing $I=I_{a,\gamma}(P)$,
\begin{align}
 \#\{r_i\in\Bcl(a,\gamma)\}&=\deg I
 =\max\{j:v(b_j)+j\gamma=g_P(a,\gamma)\},\label{eq:closedcount}\\
 \#\{r_i\in\Bop(a,\gamma)\}&=\ord_X I
 =\min\{j:v(b_j)+j\gamma=g_P(a,\gamma)\}.\label{eq:opencount}
\end{align}
For each $c\in\C$, the multiplicity of $c$ as a root of $I$ is exactly
the number of roots of $P$ in the residue subdisk
\begin{equation}\label{eq:residuedisk}
 a+t^\gamma(c+\mm).
\end{equation}
All these counts include multiplicity.
\end{theorem}
\begin{proof}
Apply \cref{lem:gaussproduct} to the linear factors. A factor with
$\lambda_i<\gamma$ reduces, after normalization by $t^{-\lambda_i}$,
to the nonzero constant $-\lc(r_i-a)$. A factor with
$\lambda_i=\gamma$ reduces to $X-\st((r_i-a)/t^\gamma)$.
A factor with $\lambda_i>\gamma$ reduces to $X$. Their Gauss values are
respectively $\lambda_i,\gamma,\gamma$. Multiplying proves
\eqref{eq:gaussroots} and \eqref{eq:initialroots}. The number of
nonconstant factors is the closed-disk count, and the number of factors
$X$ is the open-disk count. The nonzero coefficients of $I$ occur exactly
at the minimizing indices in \eqref{eq:gaussvalue}, proving the extremal
index descriptions. Finally, a scaled root has standard part $c$ exactly
when it belongs to \eqref{eq:residuedisk}.
\end{proof}

\subsection{Newton polygons without a rank-one hypothesis}
Consider $a=0$ and first remove a factor $z^m$ if $P(0)=0$. The remaining
coefficient points are $(j,v(b_j))$. Their finite lower Newton polygon
can be defined using supporting affine functions, with rational division
in $\Gamma$: the slope joining two points is
\[
 \frac{v(b_k)-v(b_j)}{k-j}\in\Gamma.
\]
This definition does not require embedding $\Gamma$ into $\R$ or drawing
an ordinary real graph. An edge of slope $-\gamma$ means that the
minimum of $v(b_j)+j\gamma$ is attained along that edge.

\begin{corollary}[Newton-polygon theorem]\label{cor:newtonpolygon}
The nonzero roots of $P$ with valuation $\gamma$ are counted by the
horizontal length of the lower edge of slope $-\gamma$. Equivalently,
that count is
\[
 \max\operatorname*{argmin}_j(v(b_j)+j\gamma)
 -\min\operatorname*{argmin}_j(v(b_j)+j\gamma).
\]
The residue polynomial of that edge determines their leading complex
coefficients, with multiplicity.
\end{corollary}
\begin{proof}
Subtract \eqref{eq:opencount} from \eqref{eq:closedcount}; then apply
\eqref{eq:initialroots}. The extremal minimizing indices are exactly
the endpoints of the supporting lower edge.
\end{proof}
Thus the one-variable tropical condition that a minimum is attained at
least twice is sufficient as well as necessary for a nonzero root at
that valuation. Its sufficiency here is a consequence of algebraic
closedness together with the exact count, not an additional convergence
assertion about a Newton algorithm.

\begin{theorem}[Valuative Rouch\'e theorem]\label{thm:valrouche}
Let $P\ne0$, let $E\in K[z]$, and suppose
\begin{equation}\label{eq:valrouchehyp}
 g_E(a,\gamma)>g_P(a,\gamma),
\end{equation}
where $g_0=+\infty$. Then $P+E$ has the same initial polynomial as $P$
at $(a,\gamma)$. It has the same root counts in the open disk, the closed
disk, and every residue subdisk \eqref{eq:residuedisk}.
\end{theorem}
\begin{proof}
After multiplying by $t^{-g_P}$ and substituting $a+t^\gamma X$, every
coefficient of $E$ has positive valuation. Thus it changes no coefficient
of the reduction. Apply \cref{thm:initialcount}.
\end{proof}
The degrees of $P$ and $P+E$ need not agree. Only the specified disk counts
are asserted. This is exactly the distinction illustrated by
\cref{ex:escape}. Unlike \cref{thm:orderedrouche}, this proof does not use
real-closed-field transfer or any version of contour integration.

\section{Coprime factor lifting with strong support control}\label{sec:hensel}
The root-count theorem locates clusters. The following factorization
constructs them when their reductions are separated. It gives a useful
polynomial counterpart of the preparation recursions in the supplied
manuscripts.

\begin{theorem}[Support-controlled coprime Hensel factorization]\label{thm:hensel}
Let $P\in\OO_K[z]$ be monic of degree $n$, and suppose its reduction is
\[
 P_0=\prod_{\nu=1}^q f_\nu,
\]
where the $f_\nu\in\C[z]$ are pairwise coprime monic polynomials of
positive degrees $d_\nu$. Write $P=P_0+E$ and let $S\subset\Gamma_{>0}$
be the union of the supports of the coefficients of $E$.
There are unique monic factors
\begin{equation}\label{eq:henselfactors}
 P=\prod_{\nu=1}^q F_\nu,\qquad
 F_\nu=f_\nu+G_\nu,\quad \deg G_\nu<d_\nu,
 \quad \supp G_\nu\subseteq S^*\setminus\{0\}.
\end{equation}
Uniqueness holds among all integral monic factors with these reductions,
not merely among factors assumed in advance to have the displayed support.
\end{theorem}
\begin{proof}
Consider the complex-linear map
\begin{equation}\label{eq:hensellinear}
 L:\bigoplus_{\nu=1}^q\C[z]_{<d_\nu}\longrightarrow\C[z]_{<n},
 \qquad (h_\nu)\longmapsto
 \sum_{\nu=1}^q h_\nu\prod_{\mu\ne\nu}f_\mu.
\end{equation}
It is injective: reducing a zero image modulo $f_\nu$, coprimality shows
that $h_\nu\equiv0\pmod{f_\nu}$, and its degree then makes it zero.
Source and target both have dimension $n$, so $L$ is an isomorphism.

Construct the coefficient polynomials of $G_\nu$ by recursion through the
well-ordered set $S^*\setminus\{0\}$. At exponent $\eta$, expand
$\prod(f_\nu+G_\nu)$. The terms containing precisely one correction
at exponent $\eta$ form $L((G_{\nu,\eta})_\nu)$. Every other nonconstant
term contains at least two positive corrections, so all its component
exponents are strictly smaller than $\eta$. Define
\[
 (G_{\nu,\eta})_\nu
 =L^{-1}\left(E_\eta-
 [t^\eta]\!\sum_{\substack{J\subseteq\{1,\ldots,q\}\\|J|\geq2}}
 \prod_{\nu\in J}G_\nu\prod_{\mu\notin J}f_\mu\right).
\]
The support lemma ensures that each coefficient uses only finitely many
products. The right side has degree less than $n$, because at least one
factor in every correction product has degree strictly below its reference
degree. The constructed factors therefore have the required degree and
support, and their product equals $P$ coefficientwise.

For uniqueness, take the least exponent where two collections of factors
differ. The terms with at least two positive corrections cannot contribute
a difference there, because they would involve a lower-index difference.
The differences at this exponent are consequently in the kernel of $L$,
which is zero. The union of the finitely many candidate supports is well
ordered, so this least-exponent argument applies without any prescribed
support bound on the candidates.
\end{proof}

The proof is the familiar coprime lifting mechanism, with its summation
requirements exposed. It applies to positive supports with infinitely many
exponents in a bounded interval and to arbitrary-rank value groups.
Repeated multiplication of corrections is legitimate by finite-word support
control; it is not justified merely by saying that a sequence of valuations
increases.

\begin{corollary}[Simple-root lifting]\label{cor:simplelift}
Every simple root $c$ of $P_0$ lifts to a unique root $c+\eta$ of $P$ with
$\eta\in\mm$. Its correction has support in $S^*\setminus\{0\}$.
When $S$ is nonempty, at its least exponent $\sigma=\min S$ the coefficient is
\[
 [t^\sigma]\eta=-\frac{E_\sigma(c)}{P_0'(c)}.
\]
The coefficient may vanish; no nonvanishing of this first correction is
assumed.
\end{corollary}
\begin{proof}
Use the coprime factors $z-c$ and $P_0/(z-c)$ in \cref{thm:hensel},
omitting the latter when it is the constant one.
The lifted degree-one factor is $z-c-\eta$. Substitution in $P$ at the
least positive exponent gives the displayed coefficient identity.
\end{proof}
For a multiple reduction root, a cluster factor still exists, but its
individual roots can require fractional exponents and further translated
Newton polygons. The divisible workspace supplies those exponents. The
factor theorem by itself does not assert that every root expansion has a
finite or effectively enumerable description.

\begin{example}[Lifting on a nondiscrete positive support]\label{ex:nondiscrete}
Let $\sigma_k=1-1/(k+2)$ and $A=\sum_{k\geq0}t^{\sigma_k}$.
The support is well ordered, although it has infinitely many elements
below $1$. The roots of $z^2-(1+A)$ are
\[
 \pm\left(1+\sum_{m\geq1}\binom{1/2}{m}A^m\right).
\]
The sum is strong and its support is contained in
$\{\sigma_k:k\geq0\}^*$. Formal binomial identities are valid
coefficientwise by the support lemma. Replacing $t$ by any ordinary
$q\in(0,1)$ makes the summands $q^{\sigma_k}$ tend to $q$, so even the
input series $A$ is not ordinarily convergent. Ordinary numerical
convergence is therefore not a hidden hypothesis of the lifting theorem.
\end{example}

\section{Exact Rolle counts and critical-point cluster trees}\label{sec:critical}
The residue field has characteristic zero, and hence every nonzero integer
has valuation zero. This seemingly small fact makes the following root
count much sharper than a general statement in positive residue
characteristic.

\begin{theorem}[Exact Rolle theorem in valuative disks]\label{thm:valuerolle}
Let $P\in K[z]$ be nonconstant. If a closed valuative disk contains
$m\geq1$ roots of $P$, it contains exactly $m-1$ roots of $P'$.
The same assertion holds for an open valuative disk containing $m\geq1$
roots. Counts include multiplicity. More generally, a disk containing
$m\geq k$ roots of $P$ contains exactly $m-k$ roots of $P^{(k)}$.
\end{theorem}
\begin{proof}
Put $g=g_P(a,\gamma)$ and $F(X)=t^{-g}P(a+t^\gamma X)$.
When the closed disk contains a root, $I=\st F$ has degree at least one.
Its derivative is nonzero, because the residue characteristic is zero.
The coefficients of $F'$ are integral and their reduction is $I'$.
It follows that
\begin{equation}\label{eq:initialderivative}
 g_{P'}(a,\gamma)=g_P(a,\gamma)-\gamma,
 \qquad I_{a,\gamma}(P')=I_{a,\gamma}(P)'.
\end{equation}
For a closed disk, its root count is $\deg I$, and the derivative count
is $\deg I'=\deg I-1$. For an open disk with a positive root count,
$I$ has a zero at $0$, and $\ord_X I'=\ord_X I-1$.
These are \eqref{eq:closedcount} and \eqref{eq:opencount} applied to the
derivative. Iterating proves the higher-derivative assertion, including
the last step from one root to zero roots.
\end{proof}

\begin{remark}[The positive-count hypothesis is necessary]
A disk containing no roots can contain critical points. For example,
$P=z^2-1$ has no roots in $\Bop(0,0)=\mm$, but $P'=2z$ has a root
there. No formula with the meaningless count $m-1=-1$ is intended.
\end{remark}

\begin{corollary}[Valuative Gauss--Lucas and nearest neighbors]\label{cor:nearest}
Every valuative disk containing all roots of $P$ contains all its critical
points. If $P$ is squarefree of degree at least two, and
\[
 \delta_i=\max_{j\ne i}v(r_i-r_j),
\]
then
\begin{equation}\label{eq:nearest}
 \max_{P'(w)=0}v(w-r_i)=\delta_i.
\end{equation}
Thus each root has a critical point at exactly its nearest-neighbor
valuation, and no critical point is closer in the valuation sense.
\end{corollary}
\begin{proof}
The first assertion is \cref{thm:valuerolle} with $m=\deg P$.
For the second, $\Bop(r_i,\delta_i)$ contains only $r_i$, whereas
$\Bcl(r_i,\delta_i)$ contains at least two roots. The open disk has no
critical point; the closed disk has at least one. Their difference is the
shell $v(w-r_i)=\delta_i$, proving \eqref{eq:nearest}.
\end{proof}
This belongs to the non-Archimedean root-location picture, which has
precedents such as \cite{ChoiLee}. It is not a proof of a Euclidean
root--critical-point distance conjecture: valuation distance ignores the
ordinary leading coefficient of a modulus, and its disks are different
from ordered-modulus disks.

\subsection{A finite root tree with explicit branch residues}
Assume for the moment that the roots $r_1,\ldots,r_n$ are distinct,
$n\geq2$. The full root set is the root node of a tree. For a node $C$
containing at least two roots, set
\[
 \delta_C=\min_{r\ne s\in C}v(r-s).
\]
Partition $C$ into children by the equivalence relation
$v(r-s)>\delta_C$. The ultrametric inequality makes this an equivalence
relation, and there are at least two children. Continue until the children
are singletons. This is a finite combinatorial tree; its depths belong to
$\Gamma$, not necessarily to $\R$.

Every node $C$ is exactly the set of roots in
$\Bcl(a,\delta_C)$ for $a\in C$. To see this, start with the root node
and observe that a child lies in one residue subdisk of its parent;
its internal depth is strictly larger than the parent's depth, whereas
all roots outside it are separated at the parent's depth or earlier.
This also proves the assertion recursively for deeper descendants.

Choose $a\in C$ and write the distinct normalized residues of its children
as $c_1,\ldots,c_k\in\C$, with child sizes $m_1,\ldots,m_k$ and
$m=\sum_jm_j$. Then
\begin{equation}\label{eq:clusterinitial}
 I_{a,\delta_C}(P)=c_*\prod_{j=1}^k(X-c_j)^{m_j}.
\end{equation}
Roots outside $C$ contribute only to the nonzero constant $c_*$.

\begin{theorem}[Critical points at a branching node]\label{thm:branch}
There are exactly $k-1$ critical points, counted with multiplicity, in
the parent disk but outside its $k$ child residue disks. Their normalized
leading residues are precisely the roots of
\begin{equation}\label{eq:branchpoly}
 W_C(X)=\sum_{j=1}^k m_j\prod_{\ell\ne j}(X-c_\ell).
\end{equation}
Each child residue disk contains exactly $m_j-1$ critical points.
The polynomial $W_C$ has degree $k-1$ and no root among the $c_j$.
\end{theorem}
\begin{proof}
Differentiate \eqref{eq:clusterinitial} and use
\eqref{eq:initialderivative}:
\[
 I_{a,\delta_C}(P')
 =c_*\prod_{j=1}^k(X-c_j)^{m_j-1}W_C(X).
\]
The leading coefficient of $W_C$ is $m\ne0$, and
$W_C(c_j)=m_j\prod_{\ell\ne j}(c_j-c_\ell)\ne0$.
The residue-subdisk count in \cref{thm:initialcount} now gives all the
assertions. In particular,
$(m-1)-\sum_j(m_j-1)=k-1$ critical points remain outside the children.
\end{proof}
The branch polynomial is an ordinary weighted incomplete polynomial of
the type in \cref{prop:weighted}. Hence its roots also lie in the ordinary
complex convex hull of the child residues. This gives a useful connection
between the order-geometric and valuation-geometric theories: the
valuation locates the scale, and an ordinary polynomial locates the
leading complex directions at that scale.

Repeated roots can be represented by terminal leaves carrying their
multiplicity. A terminal root of multiplicity $m$ contributes $m-1$
critical points at that same root. The same branch computation then uses
child multiplicities instead of cardinalities. The squarefree version is
convenient for the discriminant and stability theorems that follow.

\section{Discriminants, root trees, and coefficient precision}\label{sec:stability}
For a squarefree polynomial, the discriminant records the complete
pairwise separation data in one weighted sum:
\begin{equation}\label{eq:discvaluation}
 v(\Disc P)=(2n-2)v(a_n)+2\sum_{i<j}v(r_i-r_j).
\end{equation}
This is an immediate valuation of \eqref{eq:disc}. For monic $P$, put
\begin{equation}\label{eq:derivativedepth}
 d_i=v(P'(r_i))=\sum_{j\ne i}v(r_i-r_j).
\end{equation}
Then $\sum_i d_i=v(\Disc P)$.

\begin{proposition}[Discriminant as a sum over the root tree]\label{prop:disctree}
Let $P$ be monic and squarefree of degree $n\geq2$. With the tree of
\cref{sec:critical},
\begin{equation}\label{eq:disctree}
 v(\Disc P)=n(n-1)\delta_{\mathrm{root}}
 +\sum_{\substack{C\text{ internal}\\C\ne\mathrm{root}}}
 |C|(|C|-1)(\delta_C-\delta_{\mathrm{parent}(C)}).
\end{equation}
The depths can be negative, positive, or non-Archimedean elements of
$\Gamma$.
\end{proposition}
\begin{proof}
For a pair of distinct roots, its separation valuation is the depth of
its least common ancestor. Start with the root depth and add depth
increments for every proper internal node containing both roots; this
telescopes to that separation. Summing over unordered pairs counts
$\binom{|C|}{2}$ pairs for each node. Multiply by two and use
\eqref{eq:discvaluation}.
\end{proof}

\subsection{Normalizing a finite root problem}
The most transparent precision estimate is for monic polynomials with
integral roots. This entails no loss for an individual polynomial after
a change of scale. For
$P=z^n+\sum_{j<n}a_jz^j$, choose $\gamma\in\Gamma$ satisfying
\[
 \gamma\leq \min_{a_j\ne0}\frac{v(a_j)}{n-j}.
\]
Then $t^{-n\gamma}P(t^\gamma X)$ is monic with integral coefficients.
If the set of nonzero lower coefficients is empty, any $\gamma$ works.
Every root of a monic integral polynomial is integral: at a root of
negative valuation the highest-degree term has strictly smaller valuation
than every other term and cannot cancel. Translation before scaling is
also allowed. Precision assertions below refer to the coefficients in
these normalized coordinates.

\begin{theorem}[Root-tree stability and leading displacement]\label{thm:precision}
Let $P\in\OO_K[z]$ be monic, squarefree, and of degree $n\geq2$, with
roots $r_i$. Define
\[
 \delta_i=\max_{j\ne i}v(r_i-r_j),\qquad
 d_i=vP'(r_i),\qquad \tau(P)=\max_i(d_i+\delta_i).
\]
Let $Q=P+E$ be monic of the same degree, and put
$e=\min_jv([z^j]E)$, with $e=+\infty$ when $E=0$.
If
\begin{equation}\label{eq:precisionthreshold}
 e>\tau(P),
\end{equation}
then $Q$ is squarefree and has a unique labeling of its roots $\beta_i$
with $v(\beta_i-r_i)>\delta_i$. This labeling satisfies
\begin{equation}\label{eq:treepreservation}
 v(\beta_i-\beta_j)=v(r_i-r_j)\qquad(i\ne j).
\end{equation}
If $E(r_i)=0$, then $\beta_i=r_i$. Otherwise,
\begin{align}
 v(\beta_i-r_i)&=v(E(r_i))-d_i\geq e-d_i>\delta_i,
 \label{eq:rootprecision}\\
 \beta_i-r_i&=-\frac{E(r_i)}{P'(r_i)}(1+\varepsilon_i),
 \qquad \varepsilon_i\in\mm.\label{eq:rootfirstorder}
\end{align}
In addition $vQ'(\beta_i)=d_i$ and $v(\Disc Q)=v(\Disc P)$.
\end{theorem}
\begin{proof}
The case $E=0$ is immediate, so assume $e<+\infty$.
All $r_i$ are integral, so every $\delta_i$ and $d_i$ is nonnegative.
Choose
$\delta_i<\gamma_i<e-d_i$, possible by divisibility; the midpoint
suffices. The product formula \eqref{eq:gaussroots} gives
\[
 g_P(r_i,\gamma_i)=\gamma_i+
 \sum_{j\ne i}v(r_i-r_j)=\gamma_i+d_i<e.
\]
Every coefficient of $E(r_i+t^{\gamma_i}X)$ has valuation at least $e$:
expand each of its finitely many monomials and use the integrality of
$r_i$, the nonnegativity of $\gamma_i$, and $v\binom{j}{k}=0$.
Hence $g_E(r_i,\gamma_i)\geq e$. Valuative Rouch\'e gives one root of
$Q$ in each of the disks $\Bcl(r_i,\gamma_i)$, counted with
multiplicity. These disks are disjoint since each radius valuation is
strictly greater than every separation involving its center. Their $n$
roots exhaust the degree of $Q$, and all are simple. Denote them by
$\beta_i$.

Set $h=\beta_i-r_i$ and
$U(z)=\prod_{j\ne i}(z-r_j)$. We have $v(h)>\delta_i$ and
\[
 U(r_i+h)=P'(r_i)
 \prod_{j\ne i}\left(1+\frac{h}{r_i-r_j}\right)
 \in P'(r_i)(1+\mm).
\]
There is a polynomial $V(h)$ with
$E(r_i+h)=E(r_i)+hV(h)$ and $v(V(h))\geq e$, again by the finite
Taylor expansion at integral arguments. The root equation is therefore
\begin{equation}\label{eq:displacementequation}
 h\bigl(U(r_i+h)+V(h)\bigr)=-E(r_i).
\end{equation}
Since $e>d_i$, the bracket belongs to $P'(r_i)(1+\mm)$ and is nonzero.
If $E(r_i)=0$, this forces $h=0$. Otherwise taking its valuation and
inverting its principal unit proves \eqref{eq:rootprecision} and
\eqref{eq:rootfirstorder}.

The corrections to $r_i-r_j$ in $\beta_i-\beta_j$ have valuations
strictly greater than $v(r_i-r_j)$. The strong triangle equality proves
\eqref{eq:treepreservation}. It also shows that no other $\beta_j$ can
belong to $\Bop(r_i,\delta_i)$, proving uniqueness of the labeling.
Finally,
\[
 Q'(r_i+h)=U(r_i+h)+hU'(r_i+h)+E'(r_i+h).
\]
Here $hU'/U=\sum_{j\ne i}h/(r_i+h-r_j)$ is infinitesimal, and the last
term has valuation at least $e>d_i$. Thus $vQ'(\beta_i)=d_i$.
The discriminant identity or \eqref{eq:treepreservation} proves the last
assertion.
\end{proof}

\begin{corollary}[A coarser discriminant-only certificate]\label{cor:disccertificate}
Under the hypotheses on $P$, the condition $e>v(\Disc P)$ implies all
the conclusions of \cref{thm:precision}.
\end{corollary}
\begin{proof}
All pairwise separation valuations are nonnegative. Thus
$\delta_i\leq d_i$ and
\[
 v(\Disc P)=2\sum_{j<k}v(r_j-r_k)\geq2d_i\geq d_i+\delta_i.
\]
Take the maximum over $i$.
\end{proof}
The tree-dependent threshold can be appreciably smaller than this
certificate. It is also no larger than $2\max_i d_i$, the threshold
suggested by applying a crude scalar Hensel inequality at every root.
The theorem supplies not just existence but the exact valuation and the
leading coefficient of each nonzero displacement.

\begin{example}[The strict threshold cannot generally be weakened]\label{ex:sharp}
For $\delta>0$, take $P=z(z-t^\delta)$. Then
$d_i=\delta_i=\delta$ and $\tau(P)=2\delta$. The perturbation
\[
 Q=P+\frac14t^{2\delta}=\left(z-\frac12t^\delta\right)^2
\]
has coefficient precision exactly $\tau(P)$, but two simple roots have
collided. Thus equality at the stated threshold does not suffice in
general. This example proves sharpness of the strict boundary in this
family, not optimality of the bound for every individual perturbation.
\end{example}

\section{Univariate residue duality without contours}\label{sec:uniresidue}
For a monic polynomial $P\in K[z]$ of degree $n\geq1$, let
$A_P=K[z]/(P)$ and define
\begin{equation}\label{eq:residuetop}
 \lambda_P(H)=[z^{n-1}]\NF_P(H).
\end{equation}
Unlike the trace functional, this functional detects nilpotent directions.
It is defined by finite polynomial division, even at repeated roots and
at infinite surcomplex roots.

\begin{theorem}[Perfect residue pairing and the trace identity]\label{thm:uniduality}
The pairing $(G,H)\mapsto\lambda_P(GH)$ on $A_P$ is perfect. Its Gram
matrix in the basis $1,z,\ldots,z^{n-1}$ has determinant
$(-1)^{n(n-1)/2}$. For every polynomial $H$,
\begin{align}
 \lambda_P(H)&=\sum_{P(r)=0}\Res_{z=r}\frac{H(z)}{P(z)}\,dz,
 \label{eq:algebraicresidue}\\
 \Tr(m_H)&=\lambda_P(P'H)=\sum_{P(r)=0}m_rH(r),
 \label{eq:algebraictrace}
\end{align}
where each distinct root occurs once in these sums and $m_r$ is its
multiplicity. If $P$ is squarefree, then
\begin{equation}\label{eq:simpleresidue}
 \lambda_P(H)=\sum_{i=1}^n\frac{H(r_i)}{P'(r_i)}.
\end{equation}
\end{theorem}
\begin{proof}
For basis indices $0\leq i,j<n$, the Gram entry is zero when
$i+j<n-1$, and is one when $i+j=n-1$. Reverse the order of one basis:
the resulting matrix is triangular with diagonal entries one. Reversing
$n$ indices has sign $(-1)^{n(n-1)/2}$, proving perfection and the
determinant formula. This proof works over any commutative coefficient
ring as long as $P$ is monic.

Define the residue at $r$ as the coefficient of $(z-r)^{-1}$ in the
formal local Laurent expansion of the rational function. At a root of
multiplicity $m$, factor $P=(z-r)^mU$ with $U(r)\ne0$. The residue of
$H/P$ is the coefficient of $(z-r)^{m-1}$ in the formal Taylor expansion
of $H/U$, a finite jet. Partial fractions show that the sum of the
finite residues is the coefficient of $z^{-1}$ in the Laurent expansion
at infinity of $H/P$. The polynomial quotient of $H$ by $P$ contributes
no such coefficient, and for its remainder of degree less than $n$ that
coefficient is its $z^{n-1}$ coefficient. This proves
\eqref{eq:algebraicresidue}.

The logarithmic derivative has the finite identity
$P'/P=\sum_r m_r/(z-r)$. Its product with $H$ has residue $m_rH(r)$
at $r$, giving the first equality in \eqref{eq:algebraictrace} with the
root sum. The trace has the same root sum by the local-algebra argument
used for \eqref{eq:resultantnorm}. A simple pole at $r_i$ has residue
$H(r_i)/P'(r_i)$, proving \eqref{eq:simpleresidue}.
\end{proof}

The residue at infinity of the differential $H(z)\,dz/P(z)$ is the
negative of the coefficient just used, since $z=1/u$ gives
$dz=-u^{-2}du$. Thus the sum of all residues on the algebraic projective
line is zero. This is an algebraic residue theorem; it neither chooses
nor presupposes a surcomplex contour.

\subsection{The B\'ezoutian is a duality kernel}
Introduce two copies $X,Y$ of the coordinate. The polynomial
\begin{equation}\label{eq:unibezout}
 \Delta_P(X,Y)=\frac{P(X)-P(Y)}{X-Y}
\end{equation}
defines an element of $A_P\otimes_K A_P$. It satisfies
\[
 (X-Y)\Delta_P=0,\qquad
 (\lambda_P\otimes\id)(\Delta_P)=1.
\]
The second identity holds because the coefficient of $X^{n-1}$ is one
and the degree in $X$ is $n-1$.

For clarity, the following elementary kernel argument will also be used
in several variables. A perfect pairing identifies $A_P\otimes A_P$
with $\operatorname{End}_K(A_P)$ by
\begin{equation}\label{eq:kernelmap}
 \sum_su_s\otimes v_s\longmapsto
 \left(h\longmapsto\sum_s\lambda_P(hu_s)v_s\right).
\end{equation}
The first identity says that the operator associated with $\Delta_P$
commutes with multiplication by $z$, hence with multiplication by every
polynomial. The second says it maps $1$ to $1$. It is therefore the
identity operator. If $(b_j)$ and $(b^j)$ are residue-dual bases, this
proves
\[
 \Delta_P=\sum_jb_j(X)b^j(Y),\qquad
 P'(z)=\sum_jb_j(z)b^j(z)\quad\text{in }A_P.
\]
Taking the diagonal coefficients of multiplication by $H$ yields
$\Tr(m_H)=\lambda_P(H\sum_jb_jb^j)$, another proof of the trace identity.
The derivative is the diagonal value of the duality kernel.

\begin{corollary}[Trace discriminant versus residue determinant]\label{cor:tracegram}
In the monomial basis, the determinant of the trace-pairing matrix
$\bigl(\Tr(m_{z^{i+j}})\bigr)$ is $\Disc P$, whereas the residue-pairing
determinant is the constant $(-1)^{n(n-1)/2}$.
\end{corollary}
\begin{proof}
By \eqref{eq:algebraictrace}, the trace Gram matrix is the residue Gram
matrix times the matrix of multiplication by $P'$. Its determinant is
$(-1)^{n(n-1)/2}\Res(P,P')=\Disc P$ for monic $P$.
\end{proof}
This gives a precise collision distinction. At a repeated root the trace
pairing loses nilpotent information, whereas the residue pairing remains
perfect. Individual simple-root weights may diverge under a collision;
the finite remainder functional does not become singular.

\section{Finite polynomial systems and algebraic geometry}\label{sec:systems}
Let $F_1,\ldots,F_s\in\SC[x_1,\ldots,x_n]$ be a finite polynomial list.
Choose a divisible Hahn workspace $K$ containing their coefficients and
put $I=(F_1,\ldots,F_s)\subset K[x]$. The elementary dictionary of
affine algebraic geometry applies to this set-sized ring. We state the
basic transfer explicitly so that a class-field notation is not mistaken
for an unproved new algebraic geometry.

\begin{theorem}[Nullstellensatz in a surcomplex workspace]\label{thm:nullstellensatz}
A proper ideal $I\subset K[x]$ has a common zero in $K^n$. If
$H\in K[x]$ vanishes at every common zero, then $H^N\in I$ for some
positive integer $N$. In particular, for the finite system above,
absence of a common zero in $\SC^n$ is equivalent to the existence of
polynomials $Q_j\in K[x]$ with $\sum_jQ_jF_j=1$.
\end{theorem}
\begin{proof}
We use the classical Zariski lemma: a field finitely generated as an
algebra over a field is a finite algebraic extension. A maximal ideal
containing $I$ has quotient a finitely generated field extension of the
algebraically closed field $K$, hence quotient $K$. The images of the
coordinates give the asserted zero. This is the usual weak
Nullstellensatz; see \cite{StacksNull} for the standard algebraic theorem.

For the strong form, suppose $H$ vanishes on the zero set and form
$(I,1-yH)\subset K[x,y]$. It cannot have a zero, so the weak form says
it is the unit ideal. In an identity expressing $1$ in that ideal,
substitute $y=H^{-1}$ in $K[x,H^{-1}]$. Clearing a common power of $H$
gives $H^N\in I$. The case $H=0$ is immediate. For the final assertion,
a proper ideal would already have a zero in $K^n\subset\SC^n$;
a unit-ideal identity conversely precludes a zero in every extension.
\end{proof}

If the quotient $B=K[x]/I$ has finite dimension $D$, it decomposes into
finitely many local Artinian algebras
\begin{equation}\label{eq:finitealgebraproduct}
 B=\prod_{\xi\in Z}B_\xi,\qquad
 D=\sum_{\xi\in Z}\dim_K B_\xi.
\end{equation}
The dimensions are the algebraic local intersection multiplicities for
this zero-dimensional scheme. For a chosen finite basis, multiplication
by the coordinates is represented by commuting $D\times D$ matrices
$W_j$. Each coordinate of a point $\xi$ is a root of the corresponding
characteristic polynomial. This proves again that all surcomplex points
of a finite quotient are already in $K^n$. Moreover,
\begin{equation}\label{eq:multivartrace}
 \Tr(m_H)=\sum_{\xi\in Z}(\dim_KB_\xi)H(\xi),\qquad
 \det(m_H)=\prod_{\xi\in Z}H(\xi)^{\dim_KB_\xi}.
\end{equation}
Indeed, multiplication by $H-H(\xi)$ on each local factor is nilpotent.

A monic Gr\"obner basis computed over $K$ remains a Gr\"obner basis after
extension to another field, since its division and critical-pair
identities remain valid. Its standard monomials, and hence finite quotient
dimensions, are unchanged. This explains how elimination can be carried
out in a fixed workspace. It does not claim an algorithm for arbitrary
surreal coefficients lacking finite descriptions or computable arithmetic.
The next section gives a particularly explicit family where a basis and
all multiplicities are known before any elimination is performed.

\section{Degree-controlled complete intersections over arbitrary rings}\label{sec:globalci}
Let $A$ be a nonzero commutative ring with identity, and let
\begin{equation}\label{eq:polysystem}
 F_i(x)=x_i^{d_i}+E_i(x)\in A[x_1,\ldots,x_n],\qquad
 d_i\geq1,\qquad \deg E_i<d_i\quad(1\leq i\leq n).
\end{equation}
Here and throughout this section degree means \emph{total} degree.
Define
\[
 \mathcal B=\{\alpha\in\N^n:0\leq\alpha_i<d_i\},\qquad
 D=\prod_i d_i,\qquad \rho=\sum_i(d_i-1),\qquad
 x^{\mathrm{top}}=\prod_i x_i^{d_i-1}.
\]
The errors can be fully coupled and their coefficients can have arbitrary
size or valuation. The hypothesis is on their polynomial degree, not on
smallness.

\begin{theorem}[Global finite-free normal forms]\label{thm:globalfree}
For \eqref{eq:polysystem}, the quotient
$B=A[x]/(F_1,\ldots,F_n)$ is free over $A$ with basis
$(x^\alpha)_{\alpha\in\mathcal B}$ and rank $D$. Every polynomial $H$
has a unique normal form in their span and a finite division identity
\begin{equation}\label{eq:globalnf}
 H=\sum_iF_iQ_i+\NF_F(H).
\end{equation}
Reduction never increases total degree, and every actual use of a
relation strictly lowers the degree of the replaced monomial. All output
coefficients are polynomial expressions, with integer coefficients, in
the input coefficients. The sequence $F_1,\ldots,F_n$ is regular in
any order. These statements commute with every ring homomorphism on
coefficients.
\end{theorem}
\begin{proof}
Use the rules $x_i^{d_i}\mapsto-E_i$. Every replacement has smaller
total degree. Applied to a monomial, all branches of reduction therefore
terminate after finitely many steps; by linearity this gives a finite
reduction for each polynomial. The monomials on which no rule applies
are exactly $x^\alpha$ for $\alpha\in\mathcal B$.

We give the compatibility argument rather than assume coefficients form
a field. An ambiguity involving the rules for $i\ne j$ occurs on a
multiple of $x_i^{d_i}x_j^{d_j}$. Reducing the two indicated factors in
one order gives a multiple of $E_iE_j$; reducing in the other gives a
multiple of $E_jE_i$, the same polynomial. Any remaining reductions take
place in smaller total degrees. Induction on total degree resolves these
ambiguities, and disjoint monomial reductions commute by linearity.
Equivalently, these are monic leading terms with relatively prime leading
monomials, whose critical-pair differences reduce to zero. No division by
a coefficient is involved. It follows that the completely reduced
polynomial is independent of the reduction choices.

Each rule subtracts a multiple of one $F_i$, giving
\eqref{eq:globalnf}. Conversely $x_i^{d_i}Q$ and $-E_iQ$ have the same
normal form, so the normal-form map annihilates $F_iQ$, and hence their
ideal. An already reduced polynomial in that ideal is consequently zero.
This proves both spanning and independence of the displayed basis. It
also proves that the coefficients are obtained by finitely many additions,
subtractions, and multiplications in $A$; they therefore commute with
all coefficient homomorphisms.

For regularity, apply the same reduction to any initial subset of the
equations. Its quotient has the standard monomials with the corresponding
exponents bounded and the remaining exponents unrestricted. Its total-degree
associated graded ring is
\[
 A[x]/(x_i^{d_i}: i\text{ in that subset}),
\]
because reduction strictly lowers degree. For the next equation $F_j$,
multiplication by its highest part $x_j^{d_j}$ is injective in this graded
ring: it merely shifts the still unrestricted exponent of $x_j$ in its
monomial basis. The highest nonzero graded term of a nonzero element
therefore remains nonzero after multiplication by $F_j$. Thus $F_j$ is
a non-zero-divisor in the preceding quotient. The final quotient is
nonzero since it is free of positive rank over nonzero $A$. The same
argument works in every ordering.
\end{proof}

\begin{corollary}[A global B\'ezout count, including infinite-scale roots]\label{cor:globalbezout}
Over an algebraically closed Hahn field $K$, the system
\eqref{eq:polysystem} has total local algebraic multiplicity $D$ in
all of $\SC^n$. If its Jacobian determinant is nonzero at every solution,
there are exactly $D$ distinct solutions. There are no additional
projective solutions at infinity in the homogenized system.
\end{corollary}
\begin{proof}
Use the rank $D$ quotient and \eqref{eq:finitealgebraproduct}; workspace
independence includes every surcomplex solution. At a solution with
invertible Jacobian, the images of the equations generate the cotangent
space of the coordinate maximal ideal. Nakayama's lemma in the finite
local algebra then makes its maximal ideal zero, so that local factor
has dimension one. Finally, at the projective hyperplane $x_0=0$, the
homogenized equations read $x_i^{d_i}=0$ for every $i$. They force all
projective coordinates to vanish, which is impossible.
\end{proof}
An infinite surcomplex coordinate is still an affine coordinate in this
statement. It is not the projective point at infinity. Thus the last
sentence does not assert that all coordinates have nonnegative valuation.

\subsection{Finite support propagation without positivity}
Suppose $A$ is a Hahn coefficient ring. Let $S$ contain the support of the
coefficients of $H$, and let $T$ be the union of the coefficient supports
of the $E_i$. Neither support need be positive. If $L=\deg H\geq0$,
then the construction gives the bound
\begin{equation}\label{eq:finitepolysupport}
 \supp\NF_F(H)\ \cup\ \bigcup_i\supp Q_i
 \ \subseteq\ S+\bigcup_{k=0}^{L}kT,
\end{equation}
where $0T=\{0\}$ and the empty polynomial has zero outputs. To prove it,
follow one monomial through the reduction: each use of an error adds
one input support letter and decreases total degree by at least one.
There can be at most $L$ such uses on a branch. The witness polynomials
$Q_i$ accumulate only the preceding multipliers and obey the same bound.
Finite unions and finite Minkowski sums of well-ordered sets are well
ordered; the coefficientwise products have finite multiplicity.

This finite-degree support mechanism differs from the infinite
positive-monoid inverse in \cite{SourceResearch}. It permits negative
valuations without sacrificing support control, but only because total
degree rules out an infinite reduction chain. Neither mechanism subsumes
the other: analytic perturbations need not have bounded degree, while
polynomial perturbations need not have positive support.

\section{Multivariate B\'ezoutians and the Jacobian trace formula}\label{sec:bezoutian}
Retain \eqref{eq:polysystem} over $A$. Define a coefficient functional
on its finite-free quotient by
\begin{equation}\label{eq:multilambda}
 \lambda_F(H)=[x^{\mathrm{top}}]\NF_F(H).
\end{equation}
We prove its duality, the exact kernel representing it, and the Jacobian
trace identity directly. B\'ezoutian and residue duality are classical
subjects; \cite{BCRS,BMP} provide broader context. The point here is an
explicit proof for this family over arbitrary coefficient rings, together
with the comparison to the supplied Hahn-coherent framework.

\begin{theorem}[A coefficient-independent perfect pairing]\label{thm:multiperfect}
The pairing $(G,H)\mapsto\lambda_F(GH)$ on $B$ is perfect over $A$.
Its Gram determinant in the rectangular basis is the sign of the
permutation $\alpha\mapsto\boldsymbol d-\boldsymbol1-\alpha$ and hence
is a constant $+1$ or $-1$, independent of every coefficient of the
errors. For $\deg H<\rho$, one has $\lambda_F(H)=0$, whereas
$\lambda_F(x^{\mathrm{top}})=1$.
\end{theorem}
\begin{proof}
Reduction cannot increase degree, so
$\lambda_F(x^\alpha x^\beta)=0$ if $|\alpha|+|\beta|<\rho$.
If the sum of the degrees is $\rho$, reduction of any nonstandard
monomial strictly lowers degree, so
\begin{equation}\label{eq:gramdegree}
 \lambda_F(x^{\alpha+\beta})=
 \begin{cases}1,&\alpha+\beta=\boldsymbol d-\boldsymbol1,\\
 0,&\text{otherwise},\end{cases}
 \qquad |\alpha|+|\beta|=\rho.
\end{equation}
Order the indices $\alpha$ by nondecreasing total degree, and order the
columns by the complementary indices
$\boldsymbol d-\boldsymbol1-\alpha$. The Gram matrix is block lower
triangular, with identity diagonal blocks: degrees below the row's
complementary level give zero, and equal degrees give
\eqref{eq:gramdegree}. Its determinant in this reordered matrix is one.
Undoing the complement permutation proves the stated determinant and
therefore perfection over $A$. The last two assertions are immediate
from the same degree argument.
\end{proof}

\subsection{The divided-difference determinant}
For two coordinate tuples $X=(X_1,\ldots,X_n)$ and
$Y=(Y_1,\ldots,Y_n)$, define the polynomial matrix
\begin{equation}\label{eq:divideddiffmatrix}
 D_{ij}(X,Y)=
 \frac{F_i(Y_1,\ldots,Y_{j-1},X_j,\ldots,X_n)
       -F_i(Y_1,\ldots,Y_j,X_{j+1},\ldots,X_n)}{X_j-Y_j}.
\end{equation}
Each numerator is divisible by its displayed denominator as a polynomial
identity over $A$. Put
\begin{equation}\label{eq:multibezout}
 \Delta_F(X,Y)=\det(D_{ij}(X,Y)),\qquad
 J_F(x)=\det\left(\frac{\partial F_i}{\partial x_j}(x)\right).
\end{equation}
We view $\Delta_F$ in
$B\otimes_A B=A[X,Y]/(F(X),F(Y))$.

\begin{theorem}[B\'ezoutian kernel and Jacobian duality]\label{thm:jacobian}
Let $(b_\alpha)=(x^\alpha)_{\alpha\in\mathcal B}$ and let
$(b^\alpha)$ be its residue-dual basis for $\lambda_F$. Then
\begin{align}
 \Delta_F(X,Y)&=\sum_{\alpha\in\mathcal B}
 b_\alpha(X)b^\alpha(Y)\quad\text{in }B\otimes_A B,
 \label{eq:casimir}\\
 J_F&=\sum_{\alpha\in\mathcal B}b_\alpha b^\alpha
 \quad\text{in }B,\label{eq:jacobianeuler}\\
 \Tr_B(m_H)&=\lambda_F(J_FH)\quad(H\in B).
 \label{eq:jacobiantrace}
\end{align}
These are identities over $A$, including nonreduced coefficient rings
and parameters at which roots collide.
\end{theorem}
\begin{proof}
The columnwise telescoping identity is
\[
 D(X,Y)(X-Y)=F(X)-F(Y).
\]
Multiply by the adjugate of $D$ and pass to $B\otimes_A B$. It follows
that
\begin{equation}\label{eq:diagonalannihilator}
 (X_j-Y_j)\Delta_F=0\qquad(1\leq j\leq n).
\end{equation}
There is no matrix inverse or determinant denominator in this argument.

The total degree of row $i$ of $D$ is at most $d_i-1$.
The highest-degree contribution of that row is diagonal, coming from
$x_i^{d_i}$. Therefore $\Delta_F$ has total degree at most $\rho$,
and its homogeneous part of degree $\rho$ is
\begin{equation}\label{eq:bezouttop}
 \prod_{i=1}^n
 \left(\sum_{k=0}^{d_i-1}X_i^{d_i-1-k}Y_i^k\right).
\end{equation}
Apply $\lambda_F$ in the $X$ variables. Terms of $X$-degree less than
$\rho$ vanish. An $X$-monomial of degree $\rho$ that is nonstandard
reduces to smaller degree and also vanishes. The remaining monomial is
$X^{\mathrm{top}}$, whose coefficient in \eqref{eq:bezouttop} is one
and has no $Y$ dependence. Hence
\begin{equation}\label{eq:bezoutnormalization}
 (\lambda_F\otimes\id)(\Delta_F)=1.
\end{equation}

By \cref{thm:multiperfect}, the kernel map \eqref{eq:kernelmap}, with
$B$ and $\lambda_F$ in place of $A_P$ and $\lambda_P$, is an
isomorphism $B\otimes_A B\cong\operatorname{End}_A(B)$.
Equation \eqref{eq:diagonalannihilator} says that the operator associated
with $\Delta_F$ commutes with multiplication by each coordinate. As the
coordinates generate $B$ as an $A$-algebra, this operator is $B$-linear.
Equation \eqref{eq:bezoutnormalization} says that it sends $1$ to $1$,
so it is the identity. Its unique representing tensor is the right side
of \eqref{eq:casimir}.

The multiplication map $B\otimes_A B\to B$ substitutes $X=Y=x$.
Each divided difference then becomes the corresponding partial derivative,
so the diagonal of $\Delta_F$ is $J_F$. This proves
\eqref{eq:jacobianeuler}. Finally the coefficient of $b_\alpha$ in
$Hb_\alpha$ is $\lambda_F(Hb_\alpha b^\alpha)$. Summing these diagonal
matrix coefficients gives \eqref{eq:jacobiantrace}.
\end{proof}

\begin{corollary}[Simple-root weights and Euler--Jacobi vanishing]\label{cor:eulerjacobi}
Over an algebraically closed Hahn field $K$, suppose every solution of
\eqref{eq:polysystem} is simple. Then
\begin{equation}\label{eq:multisimpleresidue}
 \lambda_F(H)=\sum_{\xi\in Z}\frac{H(\xi)}{J_F(\xi)}.
\end{equation}
Consequently
\begin{equation}\label{eq:eulerjacobi}
 \sum_{\xi\in Z}\frac{H(\xi)}{J_F(\xi)}=0
 \quad\text{if }\deg H<\rho,
 \qquad
 \sum_{\xi\in Z}\frac{\xi^{\mathrm{top}}}{J_F(\xi)}=1.
\end{equation}
For arbitrary multiplicities, let $e_\xi$ be the local-factor idempotent
in \eqref{eq:finitealgebraproduct}. Then
\[
 \lambda_{F,\xi}(H):=\lambda_F(e_\xi H)
\]
is a perfect local residue functional and
\begin{equation}\label{eq:localjacobianlength}
 \lambda_{F,\xi}(J_FH)=(\dim_KB_\xi)H(\xi).
\end{equation}
\end{corollary}
\begin{proof}
Different local factors are orthogonal for the residue pairing because
their products are zero. Its restriction to each factor is perfect.
For a simple factor $K e_\xi$, write
$\lambda_F(e_\xi H)=c_\xi H(\xi)$. The trace identity applied to
$e_\xi$ gives $1=c_\xi J_F(\xi)$, proving
\eqref{eq:multisimpleresidue}. Now use \cref{thm:multiperfect} for the
degree bounds. In a factor of length $\mu$, multiplication by $H$
has trace $\mu H(\xi)$, and the global identity applied to $e_\xi H$
gives \eqref{eq:localjacobianlength}.
\end{proof}
The simple-root formula is not used to define the functional. Defining it
that way would obscure the multiple-root case and falsely suggest that
division by a discriminant is necessary. The top-normal-form definition,
the perfect Gram matrix, and the B\'ezoutian all remain defined at every
collision.

\section{Hahn-coherent parameters and comparison with the supplied residues}\label{sec:comparison}
Let $U\subseteq\C^s$ be an ordinary domain and, following the supplied
manuscripts, let
\[
 \HH(U)=\left\{\sum_{\gamma\in S}h_\gamma t^\gamma:
 S\subseteq\Gamma\text{ well ordered},\quad h_\gamma\in\OO(U)\right\}.
\]
Here $\OO(U)$ denotes ordinary holomorphic functions. The subring
$\Hp(U)$ consists of nonnegative supports. At $u=c+\varepsilon\in U\sh$,
evaluation uses strong Taylor substitution
\[
 \left(\sum h_\gamma t^\gamma\right)^{\sharp}(c+\varepsilon)
 =\sum_{\gamma,\alpha}
 \frac{\partial^\alpha h_\gamma(c)}{\alpha!}
 t^\gamma\varepsilon^\alpha.
\]
The support is contained in $S+T_\varepsilon^*$, where
$T_\varepsilon$ is the finite union of the positive supports of the
displacement coordinates. This is the evaluation convention proved in
\cite{SourceAnalysis,SourceResearch}, not a new interpretation of a fine
limit.

\begin{theorem}[Global polynomial families on one ordinary parameter domain]\label{thm:polyparameters}
Suppose the finitely many coefficients in \eqref{eq:polysystem} lie in
$A=\HH(U)$ or $\Hp(U)$. Then its global polynomial quotient is finite
free of rank $D$ over $A$. Normal forms, multiplication matrices, the
residue Gram matrix and its inverse, and the B\'ezoutian trace identity
all have coefficients holomorphic on the same ordinary domain $U$.
They commute with every allowed finite coherent parameter substitution
and with specialization to a finite surcomplex parameter.
No positivity of the errors is required for the version over $\HH(U)$.
\end{theorem}
\begin{proof}
Apply \cref{thm:globalfree,thm:multiperfect,thm:jacobian} to the indicated
coefficient ring. Normal-form computations are finite and obey
\eqref{eq:finitepolysupport}; they only add and multiply functions already
holomorphic on $U$. The inverse Gram matrix is its adjugate divided by
the constant determinant $\pm1$, so it too involves only finite polynomial
operations. All other constructions are finite determinants or matrix
operations. No ordinary neighborhood is shrunk. Since evaluation and
allowed coherent parameter substitutions are ring homomorphisms, the
base-change assertion of \cref{thm:globalfree} proves compatibility.
\end{proof}

\subsection{Identifying the analytic residue in a polynomial subclass}
Now suppose every coefficient of every $E_i$ has positive Hahn support.
For this paragraph the parameter ring is integral. On an ordinary product
of circles in the fiber variables, the supplied research manuscript defines
its residue by the strong coefficient formula
\begin{equation}\label{eq:sourceLambda}
 \Lambda_F(H)=\sum_{k\in\N^n}(-1)^{|k|}
 \left[x_1^{d_1(k_1+1)-1}\cdots x_n^{d_n(k_n+1)-1}\right]
 \left(H\prod_i E_i^{k_i}\right).
\end{equation}
The formula follows by expanding $1/F_i$ about $x_i^{d_i}$ and applying
the ordinary Cauchy coefficient formula at each Hahn exponent. Positive
supports make this family strongly summable. Its values are coherent
parameter functions, and $\Lambda_F(F_iH)=0$: after cancellation of
$F_i$, the coefficientwise contour integrand is holomorphic in $x_i$
on its disk, so its $x_i$ integral is zero. These are precisely the
functional and annihilation property proved in \cite{SourceResearch}.

\begin{theorem}[Polynomial identification of the Jacobian trace element]\label{thm:analyticcomparison}
Under the lower-total-degree and positive-support hypotheses just stated,
\begin{equation}\label{eq:residuecomparison}
 \Lambda_F(H)=\lambda_F(H)\quad\text{for every polynomial }H.
\end{equation}
Consequently the trace element of the supplied residue pairing, restricted
to this polynomial family, is exactly the class of $J_F$:
\[
 \Tr(m_H)=\Lambda_F(J_FH).
\]
All solutions in a specialized surcomplex fiber are infinitesimal, so its
global polynomial count $D$ agrees with the finite-monad analytic count.
\end{theorem}
\begin{proof}
It suffices first to evaluate \eqref{eq:sourceLambda} on a standard
monomial $H=x^\alpha$, whose degree is at most $\rho$.
For $k\ne0$, the extraction degree in that formula is
$\rho+\sum_i d_i k_i$, whereas the numerator has degree at most
\[
 |\alpha|+\sum_i(d_i-1)k_i
 \leq\rho+\sum_i d_i k_i-|k|.
\]
It is strictly smaller. Thus every $k\ne0$ term is zero, and the
$k=0$ term is one precisely for the top monomial and zero for the other
basis monomials. Since both functionals annihilate the defining ideal,
normal forms extend equality to every polynomial. Their Gram matrices in
the shared rectangular basis coincide, so the unique element representing
the trace under the perfect pairing is $J_F$ by \cref{thm:jacobian}.

For the location of roots, specialize the parameters in an algebraically
closed Hahn workspace. The multiplication matrices of $x_i$ have integral
coefficients. Reducing their entries modulo $\mm$ gives the nilpotent
coordinate shifts in $K[x]/(x_1^{d_1},\ldots,x_n^{d_n})$, reduced over
$\C$. Each characteristic polynomial therefore has the form
$T^D+\sum_{j<D}c_jT^j$ with every $c_j\in\mm$. A noninfinitesimal
root would have finite inverse, and division by its $D$th power would
express zero as one plus a finite sum of infinitesimals. This is impossible.
All eigenvalues, including every solution coordinate, are infinitesimal.
The finite polynomial quotient and the analytic quotient both have the
same rectangular basis and the same defining multiplication relations;
the inclusion of polynomials therefore induces an isomorphism between
these finite quotients.
\end{proof}

This settles a precisely delimited comparison left unproved in the supplied
research manuscript: the Jacobian identity for its residue functional in
the lower-degree polynomial family. It does not settle the comparison for
arbitrary holomorphic errors, nor for polynomials with higher-degree
positive errors that can create roots outside the finite monad.
\Cref{ex:escape} is a warning against that extrapolation. The general
classical theory of complete-intersection residues already contains
Jacobian identities; the contribution here is a direct normalization and
support-compatible identification with the particular supplied functional,
not a claim that the classical identity was previously unknown.

\section{Worked examples across finite, infinite, and colliding scales}\label{sec:examples}
In the first two examples $t=\omega^{-1}$, so $v(t)=1$. The notation
$o_v(t^a)$ means an element of valuation strictly greater than $a$;
it is a valuation inequality, not a statement about a fine-topological
limit. Ordinary formal power-series $O(t^a)$ notation will only be used
when the series really has integer powers and begins at degree $a$.

\subsection{A cubic with two infinite roots and one infinitesimal root}
Consider
\begin{equation}\label{eq:threecubic}
 P(z)=z^3-t^{-2}z+t^3.
\end{equation}
Its nonzero coefficient points are $(0,3),(1,-2),(3,0)$.
The two lower slopes are $-5$ and $1$, of horizontal lengths one and
two. Thus its root valuations are $5,-1,-1$.
At the scale $\gamma=5$, the initial polynomial is $1-X$; at the scale
$\gamma=-1$, it is $X^3-X$. The nonzero residues of the latter give
the two infinite leading terms $\pm t^{-1}$, while its zero residue
contains the smaller root.

For the small root put $z=t^5u$ and $s=t^{12}$. The equation becomes
$u=1+s u^3$. Its reduction has the simple root $u=1$, so the
support-controlled lifting theorem constructs a unique solution
$u\in1+s\C[[s]]$. Successive coefficient comparison gives
\begin{equation}\label{eq:smallcubicexpansion}
 z=t^5+t^{17}+3t^{29}+12t^{41}+O(t^{53}).
\end{equation}
For the two large roots put $z=t^{-1}u$. Then
$u^3-u+t^6=0$, with simple reduction roots $u=\pm1$.
Their first corrections are both $-t^6/2$, giving
\begin{equation}\label{eq:largecubicexpansion}
 z_\pm=\pm t^{-1}-\frac12t^5+O(t^{11}).
\end{equation}
The two descriptions use different root scales but one finite polynomial.
The root count proves that no additional roots are hidden at a smaller
infinitesimal or larger infinite scale.

\subsection{A four-root tree and a sharper stability threshold}\label{subsec:treeexample}
Let
\begin{equation}\label{eq:treequartic}
 P(z)=z(z-t^3)(z-t)(z-t-t^2),
\end{equation}
with roots in the displayed order. The pair $\{0,t^3\}$ has separation
valuation $3$, the pair $\{t,t+t^2\}$ has separation valuation $2$,
and each of the four cross-pair separations has valuation $1$.
Thus the root node has depth $1$, with two children of size two at
depths $3$ and $2$.

The discriminant depth, derivative depths, and nearest-neighbor depths
are
\begin{align*}
 v(\Disc P)&=2(3+2+4)=18,\\
 (d_1,d_2,d_3,d_4)&=(5,5,4,4),\\
 (\delta_1,\delta_2,\delta_3,\delta_4)&=(3,3,2,2),
 \qquad \tau(P)=8.
\end{align*}
The tree formula gives the same discriminant value:
\[
 4\cdot3\cdot1+2\cdot1(3-1)+2\cdot1(2-1)=18.
\]
The initial polynomial at the parent scale is $X^2(X-1)^2$.
Its branch polynomial is $2(X-1)+2X=4X-2$. The three critical points
therefore have descriptions
\begin{equation}\label{eq:quarticcritical}
 \frac12t^3+o_v(t^3),\qquad
 t+\frac12t^2+o_v(t^2),\qquad
 \frac12t+o_v(t).
\end{equation}
The first two lie in the two children, one in each; the third lies at
the parent branch outside the children. These three already exhaust
the degree of $P'$.

Now set $Q=P+t^9$. Its coefficient precision $9$ exceeds $\tau(P)=8$,
although it does not exceed $v(\Disc P)=18$. The sharper theorem applies
and preserves every pairwise separation. Since $E(r_i)=t^9$ at every
root, the root displacement valuations are exactly $(4,4,5,5)$.
The leading coefficients of $P'(r_i)$ are $(-1,1,-1,1)$ at their
respective valuations. Thus the matching roots are
\begin{align*}
 \beta_1&=t^4+o_v(t^4),&
 \beta_2&=t^3-t^4+o_v(t^4),\\
 \beta_3&=t+t^5+o_v(t^5),&
 \beta_4&=t+t^2-t^5+o_v(t^5).
\end{align*}
This example separates three notions: the aggregate discriminant depth,
the depth needed to stabilize the full tree, and the individual accuracy
of each root after a specified perturbation.

\subsection{A genuinely higher-rank root expansion}\label{subsec:highrank}
Take the ordered exponent group $\Gamma=\Q\omega+\Q\subset\No$ and
consider
\begin{equation}\label{eq:highrankpoly}
 P(z)=z^2-t^\omega z+t^{2\omega+1}.
\end{equation}
The change $z=t^\omega u$ gives $u^2-u+t=0$. Its two simple reduction
roots yield
\[
 u_- =\frac{1-\sqrt{1-4t}}2
 =t+t^2+2t^3+5t^4+14t^5+\cdots,
 \qquad u_+=1-u_-.
\]
Consequently
\begin{equation}\label{eq:highrankroots}
 z_-=t^{\omega+1}+t^{\omega+2}+2t^{\omega+3}+5t^{\omega+4}+\cdots,
 \qquad z_+=t^\omega-z_-.
\end{equation}
Their valuations are $\omega+1$ and $\omega$. The binomial series is
strongly summable and solves the polynomial identity coefficientwise;
no real-valued parametrization of $\omega$ is involved.

Even the simpler strong sum $\sum_{m\geq0}t^m=(1-t)^{-1}$ illustrates
why ordinary sequential intuition is unsafe in this workspace.
The difference from its $N$th partial sum has valuation $N+1$, which
never reaches $\omega$. Hence the partial sums do not converge to the
strong sum in the valuation topology of $K_\Gamma$: the neighborhood
requiring error valuation greater than $\omega$ is never entered.
In particular, the strong sum is not a fine-topological limit either.
An increasing sequence of valuations need not be cofinal in an
arbitrary-rank ordered group.

\subsection{A quartic residue pairing through a double collision}\label{subsec:quarticcollision}
Start over the polynomial parameter ring $\Q[s]$ with
\[
 P_s(z)=(z^2-s^2)^2.
\]
For every value of $s$, the quotient has the basis $1,z,z^2,z^3$,
and $\lambda_{P_s}$ is the coefficient of $z^3$ in the remainder.
Its residue Gram determinant is $(-1)^6=1$.
When $s\ne0$ in a characteristic-zero field, the two roots $s,-s$
both have multiplicity two. The local residue formula gives, for any
polynomial $H$,
\begin{equation}\label{eq:quarticcollisionresidue}
 \lambda_{P_s}(H)=
 \frac{H'(s)+H'(-s)}{4s^2}
 +\frac{H(-s)-H(s)}{4s^3}.
\end{equation}
Indeed, at $s$ the coefficient of $(z-s)^{-1}$ in
$H/((z-s)^2(z+s)^2)$ is $H'(s)/(4s^2)-H(s)/(4s^3)$;
the computation at $-s$ gives the other two terms. In contrast,
\[
 \Tr(m_H)=2H(s)+2H(-s).
\]
At the parameter value $s=0$ there is one root of length four and
\[
 \lambda_{P_0}(H)=\frac{H'''(0)}6,\qquad \Tr(m_H)=4H(0).
\]
These are not ordinary limits required to define the functionals:
polynomial division over $\Q[s]$ defines them at all parameters,
and the equality of formulas for nonzero $s$ is a rational identity.
After cancellation, the residue expression has the indicated polynomial
specialization. One may then evaluate $s$ at an infinitesimal surreal.
There is no field homomorphism sending an invertible fixed Hahn monomial
$t$ to zero; the collision parameter specialization is made before
passing to a field in which that parameter has been inverted.

\subsection{A six-dimensional coupled algebra with an explicit dual basis}\label{subsec:sixpoint}
Let $a,b,c,d$ be arbitrary surcomplex coefficients, or indeterminates in
a commutative coefficient ring, and consider
\begin{equation}\label{eq:sixsystem}
 F_1=x^2-ay-b,\qquad F_2=y^3-cx-d.
\end{equation}
The degree hypotheses hold with $(d_1,d_2)=(2,3)$.
The quotient has dimension or rank six and basis
\begin{equation}\label{eq:sixbasis}
 (1,y,y^2,x,xy,xy^2),\qquad \lambda_F(H)=[xy^2]\NF_F(H).
\end{equation}
No nonzero condition on $a,b,c,d$ is required. For a second tuple $(u,v)$,
\[
 D(X,Y)=\begin{pmatrix}x+u&-a\\-c&y^2+yv+v^2\end{pmatrix},
 \qquad
 \Delta_F=(x+u)(y^2+yv+v^2)-ac.
\]
Reading its coefficients in the first basis gives the residue-dual basis
\begin{equation}\label{eq:sixdual}
 (xy^2-ac,xy,x,y^2,y,1).
\end{equation}
Its diagonal is
\begin{equation}\label{eq:sixjacobian}
 J_F=6xy^2-ac,
\end{equation}
and the residue Gram determinant is $-1$ for every coefficient choice.
Thus the trace formula here is
$\Tr(m_H)=\lambda_F((6xy^2-ac)H)$, including all collisions.

When $c\ne0$, ordinary elimination gives
\begin{equation}\label{eq:sixelimination}
 y^6-2dy^3-ac^2y+d^2-bc^2=0,
 \qquad x=\frac{y^3-d}{c}.
\end{equation}
This is an alternative one-variable presentation. At $c=0$ it is the
division by $c$ that fails, not the six-element basis or the residue
theory. This distinction is useful for parameter computations: a global
basis avoids unnecessary exceptional denominators introduced by a
particular elimination choice.

For a concrete infinitesimal splitting take $a=c=t>0$ and $b=d=0$.
The roots are
\begin{equation}\label{eq:sixroots}
 (0,0),\qquad
 (t^{4/5}\zeta^3,t^{3/5}\zeta),\quad \zeta^5=1.
\end{equation}
To see this, solve $x=y^3/t$ and obtain $y(y^5-t^3)=0$.
At the origin the Jacobian is $-t^2$; at each of the other five roots
it is $5t^2$. All six roots are simple and the exact residue formula is
\begin{equation}\label{eq:sixweights}
 \lambda_F(H)=-\frac{H(0,0)}{t^2}
 +\frac1{5t^2}\sum_{\zeta^5=1}
 H(t^{4/5}\zeta^3,t^{3/5}\zeta).
\end{equation}
Each individual weight is infinite, yet $\lambda_F(1)=0$ by cancellation,
and $\lambda_F(xy^2)=1$. The exact symbolic checks verify these identities
on the six-element basis; normal forms then establish them for every
polynomial. In the underlying parameter family the value $a=c=0$,
$b=d=0$ gives the length-six origin with algebra
$K[x,y]/(x^2,y^3)$ and the same nondegenerate residue determinant.

\section{Computational meaning, limitations, and further directions}\label{sec:limits}
\subsection{What is finite and what may be transfinite}
Factorization exists in a set-sized algebraically closed workspace, but
existence does not give a finite input representation for every surreal
coefficient. There are three different computational statements here.
Finite polynomial identities, matrix traces, normal forms for
\eqref{eq:polysystem}, and determinant formulas are executable with exact
arithmetic whenever coefficient arithmetic is executable. Initial-polynomial
root counts additionally require accessible coefficient valuations and
leading coefficients. Strong Hensel recursions require a usable description
of supports and their finite word decompositions at the exponents to be
computed. An arbitrary well-ordered surreal support need not supply such
an effective description.

Likewise, a finite root cluster tree records only finitely many pairwise
separation values. It does not encode all transfinite terms of the
individual roots. Published work on lengths of algebraic roots in Hahn
fields, such as \cite{KnightLange}, addresses a different and more detailed
support-complexity problem. We do not infer a uniform finite root-expansion
algorithm from the finite branching tree.

\subsection{Hypotheses that cannot be silently removed}
The exact disk Rolle theorem uses residue characteristic zero. Over a
field of positive residue characteristic, differentiation can erase a
leading coefficient and change the initial derivative polynomial.
The ordered Gauss--Lucas theorem uses a real closed ordered real part
and its conjugation; algebraic closedness alone does not define its convex
hull. The precision theorem uses integral normalized roots and a strict
inequality; \cref{ex:sharp} shows an actual boundary collision.
The global complete-intersection theorem uses strict lower total degree;
positive support without that restriction does not prevent global roots
from escaping the finite monad, as \cref{ex:escape} shows.

The article does not identify arbitrary fine-locally analytic functions
with polynomials, does not extend the field valuation to a real-valued
absolute value without an additional rank-one assumption, and does not
make a field specialization $t\mapsto0$. Residues here are either finite
formal Laurent coefficients or the specifically identified coherent
contour functional; those interpretations are never conflated with
integration along fine-continuous real-parameter loops.

\subsection{Directions justified by the proved results}
One extension is to replace the coordinate-power highest terms by a
homogeneous regular sequence whose quotient has a known finite basis.
A suitable graded division and duality normalization could then replace
the rectangular argument. The existing general theory of complete
intersections and B\'ezoutians is relevant, but a fixed-basis arbitrary-ring
statement must specify its hypotheses rather than assume freeness.

A second direction is perturbation theory for repeated-root trees.
The initial-polynomial theorem already preserves residue-cluster counts,
but a theorem preserving individual labels is impossible through genuine
collisions. A natural replacement would preserve a nested collection of
cluster algebras with explicit coefficient-precision thresholds.

A third direction is the analytic residue comparison for the full class
of positive-support holomorphic perturbations in \cite{SourceResearch}.
Our proof identifies the Jacobian trace element for lower-degree
polynomial errors because a finite total-degree bound normalizes the
B\'ezoutian and the contour functional. The general analytic comparison
needs another argument. This is a boundary of the proof supplied here,
not a claim that the relevant classical compatibility problem is an
unresolved published conjecture.

\section{Theorem provenance and conclusion}\label{sec:provenance}
The following audit separates established principles from the extensions
of the attached framework. A direct proof in this article is not itself
evidence of mathematical priority.

\begin{longtable}{@{}L{0.26\textwidth}L{0.64\textwidth}@{}}
\toprule
Result & Provenance and precise role here\\
\midrule
\endhead
Workspace algebraic closure & Established surreal/Hahn theory
\cite{Conway,Gonshor,vdDE,Poonen}; used to reduce every finite task to
ordinary set-sized algebra.\\[0.45em]
Factorization, CRT, resultants, Nullstellensatz & Classical algebra over
an algebraically closed field; the workspace interpretation prevents
improper class-size assertions.\\[0.45em]
Gauss--Lucas, Hermite forms, ordered Rouch\'e & Classical finite algebraic
geometry or real-closed-field transfer. Direct proofs and explicit
first-order scope are provided.\\[0.45em]
Initial polynomial and Hensel lifting & Classical valued-field mechanisms;
the statements here retain arbitrary-rank surreal exponents and strong
support bookkeeping.\\[0.45em]
Disk Rolle and branch residues & Consequences of exact initial-polynomial
differentiation in residue characteristic zero, organized as a finite
critical-point cluster theory. Non-Archimedean root-location precedents
are acknowledged \cite{ChoiLee}.\\[0.45em]
Tree discriminants and root precision & Explicit deductions from pairwise
factorization and valuative root counting; the proof gives a tree-dependent
threshold and the exact leading displacement. No priority claim is made.\\[0.45em]
Rectangular global quotient & A monic degree-decreasing polynomial
construction over arbitrary rings, rather than the analytic positive-support
inverse in the supplied manuscript.\\[0.45em]
B\'ezoutian and Jacobian duality & Classical principles with broad
predecessors \cite{BCRS,BMP}; a direct arbitrary-ring proof is given for
the stated degree-controlled family.\\[0.45em]
Comparison with supplied contour residues & A proved identification for
positive lower-degree polynomial errors. It establishes the Jacobian
trace element in that subclass without asserting the full analytic case.\\
\bottomrule
\end{longtable}

Primary-source bibliographic checks were performed on September 21, 2026.
The literature discussion is selective, not an exhaustive search for all
equivalent formulations. In particular, the word ``surcomplex'' does not
create a new theorem when its proof only uses algebraic closedness.

The resulting polynomial algebra is nevertheless richer than merely
repeating the fundamental theorem of algebra. Ordered geometry controls
convex root location, the valuation resolves arbitrary scales and cluster
branching, and finite quotient algebras preserve multiplicity and residue
information through collisions. The interaction is concrete: a branch
polynomial gives the ordinary leading directions of critical points at a
surreal scale; a discriminant is a weighted tree length; a coefficient
threshold fixes every root separation; and a polynomial B\'ezoutian gives
a collision-stable Jacobian trace formula in the same parameter category
as the supplied analytic work. These are compatible pieces of one finite,
explicitly scoped surcomplex polynomial theory.

\appendix
\section{Notation and operational reference}\label{app:notation}
\begin{longtable}{@{}L{0.30\textwidth}L{0.60\textwidth}@{}}
\toprule
Notation & Meaning\\
\midrule
\endhead
$\SC=\No[i]$ & Surcomplex class field; every actual polynomial calculation
is made in a set-sized workspace.\\[0.35em]
$K_\Gamma=\C((t^\Gamma))$ & Algebraically closed Hahn workspace for
divisible $\Gamma\subseteq\No$.\\[0.35em]
$t^\gamma=\omega^{-\gamma}$ & Conway monomial; $v(t^\gamma)=\gamma$.\\[0.35em]
$\OO_K,\mm,\st$ & Finite valuation ring, infinitesimal ideal, and
standard-part map to $\C$.\\[0.35em]
$\Bcl(a,\gamma),\Bop(a,\gamma)$ & Points with $v(z-a)\geq\gamma$ and
$v(z-a)>\gamma$, respectively.\\[0.35em]
$g_P(a,\gamma),I_{a,\gamma}(P)$ & Weighted coefficient minimum and its
normalized complex initial polynomial.\\[0.35em]
$\delta_C,\delta_i,d_i$ & Cluster depth; nearest-neighbor depth; derivative
valuation at a simple root. The symbol $d_i$ instead denotes an integer
coordinate degree in the complete-intersection sections.\\[0.35em]
$\NF_P,\NF_F$ & Unique polynomial remainder or rectangular normal form.\\[0.35em]
$\lambda_P,\lambda_F$ & Coefficient of the top basis monomial in the
normal form; the algebraic residue functional.\\[0.35em]
$\Delta_F,J_F$ & Divided-difference B\'ezoutian and Jacobian determinant.\\[0.35em]
$D,\rho$ & Rank $\prod_i d_i$ and top degree $\sum_i(d_i-1)$ of the
rectangular quotient.\\[0.35em]
$\Lambda_F$ & The supplied coefficientwise coherent contour residue,
identified with $\lambda_F$ under \cref{thm:analyticcomparison}.\\[0.35em]
$o_v(t^a)$ & An error of valuation strictly greater than $a$, with no
fine-limit assertion.\\
\bottomrule
\end{longtable}

\section{Exact verification artifacts and reproducibility}\label{app:verification}
The archive contains \texttt{verify\_examples.py} and the report produced
by running it. The executed report records 116 passing checks using Python 3.13.5 and
SymPy 1.14.0. The script uses exact rational and symbolic arithmetic,
not floating-point approximations. Its checks cover the displayed cubic
expansions, the quartic separation and discriminant data, the leading
root corrections, the quadratic boundary collision, the higher-rank
example's ordinary coefficient recursion, the quartic residue formula,
and finite multivariate quotient examples. It tests the B\'ezoutian dual
matrix, the residue Gram determinant, multiplication relations and trace
identities, including the symbolic six-dimensional family.

These checks verify finite instances and formulas. They do not verify
arbitrary supports, class-set foundations, the transfinite Hensel recursion,
real-closed-field transfer, or the universal statements over all coefficient
rings. Those claims rest on the proofs, which have not been independently
refereed or machine-checked.

The \LaTeX{} source is standalone and includes the bibliography. Compile
with
\begin{verbatim}
latexmk -pdf surcomplex_polynomial_algebra.tex
\end{verbatim}
or run \texttt{pdflatex} repeatedly until all references stabilize.
Run the checks with Python and SymPy installed:
\begin{verbatim}
python verify_examples.py
\end{verbatim}
No external image files or bibliography processor are required.

\clearpage
\begin{thebibliography}{99}
\small
\addcontentsline{toc}{section}{References}

\bibitem{SourceAnalysis}
\emph{Surcomplex Analysis: Infinitesimal Calculus, Hahn-Coherent
Holomorphy, and Contour Theory}.
User-supplied unpublished manuscript, dated September 21, 2026,
file \texttt{surcomplex\_analysis(3).tex}.
The source of the strong-summation, workspace, and coherence conventions
used here; not an independently verified publication.

\bibitem{SourceResearch}
\emph{Hartogs Coherence, Finite Maps, and Residue Duality over the
Surcomplex Numbers: Support-controlled several-variable extensions of
Hahn-coherent analysis}.
User-supplied unpublished manuscript, dated September 21, 2026,
file \texttt{surcomplex\_research.tex}.
Especially its finite-free normal forms, multidimensional residue
functional, relative residue duality, and trace-element discussion.

\bibitem{Conway}
John H. Conway, \emph{On Numbers and Games}, second edition,
A K Peters, 2001.
\href{https://www.routledge.com/On-Numbers-and-Games/Conway/p/book/9781568811277}{Publisher record}.

\bibitem{Gonshor}
Harry Gonshor, \emph{An Introduction to the Theory of Surreal Numbers},
London Mathematical Society Lecture Note Series \textbf{110},
Cambridge University Press, 1986.
\href{https://doi.org/10.1017/CBO9780511629143}{doi:10.1017/CBO9780511629143}.

\bibitem{vdDE}
Lou van den Dries and Philip Ehrlich,
``Fields of surreal numbers and exponentiation,''
\emph{Fundamenta Mathematicae} \textbf{167} (2001), 173--188.
\href{https://doi.org/10.4064/fm167-2-3}{doi:10.4064/fm167-2-3}.

\bibitem{Poonen}
Bjorn Poonen, ``Maximally complete fields,''
\emph{L'Enseignement Math\'ematique} (2) \textbf{39} (1993), 87--106.
\href{https://math.mit.edu/~poonen/papers/amsval.pdf}{Author-hosted text}.
The algebraic-closure input is Corollary 4.

\bibitem{Swan}
Richard G. Swan, \emph{Tarski's Principle and the Elimination of Quantifiers},
author's expository notes.
\href{https://www.math.uchicago.edu/~swan/expo/Tarski.pdf}{Author-hosted text}.

\bibitem{Eisermann}
Michael Eisermann,
``The Fundamental Theorem of Algebra made effective: an elementary
real-algebraic proof via Sturm chains,''
\href{https://arxiv.org/abs/0808.0097}{arXiv:0808.0097}.
Cited for real-algebraic complex root counting rather than a fine-topological
contour interpretation.

\bibitem{ChoiLee}
Daebeom Choi and Seewoo Lee, ``Non-archimedean Sendov's Conjecture,''
\emph{p-Adic Numbers, Ultrametric Analysis and Applications}
\textbf{14} (2022), 77--80.
\href{https://doi.org/10.1134/S2070046622010058}{doi:10.1134/S2070046622010058};
\href{https://arxiv.org/abs/2106.11155}{arXiv:2106.11155}.

\bibitem{Zhang}
Teng Zhang, ``When D-companion matrix meets incomplete polynomials,''
\emph{Journal of Mathematical Analysis and Applications}
\textbf{539} (2024), article 128466.
\href{https://doi.org/10.1016/j.jmaa.2024.128466}{doi:10.1016/j.jmaa.2024.128466};
\href{https://arxiv.org/abs/2310.08930}{arXiv:2310.08930}.

\bibitem{BCRS}
Eberhard Becker, Jean-Paul Cardinal, Marie-Fran\c{c}oise Roy,
and Zbigniew Szafraniec,
``Multivariate Bezoutians, Kronecker symbol and Eisenbud--Levine formula,''
in \emph{Algorithms in Algebraic Geometry and Applications},
Progress in Mathematics \textbf{143}, Birkh\"auser, 1996, 79--104.
\href{https://doi.org/10.1007/978-3-0348-9104-2_5}{doi:10.1007/978-3-0348-9104-2\_5}.

\bibitem{BMP}
Thomas Brazelton, Stephen McKean, and Sabrina Pauli,
``B\'ezoutians and the $\mathbb A^1$-degree,''
\emph{Algebra \& Number Theory} \textbf{17} (2023), 1985--2012.
\href{https://doi.org/10.2140/ant.2023.17.1985}{doi:10.2140/ant.2023.17.1985};
\href{https://arxiv.org/abs/2103.16614}{arXiv:2103.16614}.

\bibitem{StacksNull}
The Stacks Project Authors,
``Hilbert Nullstellensatz,'' Theorem 10.34.1, Tag 00FV.
\href{https://stacks.math.columbia.edu/tag/00FV}{Stacks Project, Tag 00FV}.
Accessed September 21, 2026.

\bibitem{KnightLange}
Julia F. Knight and Karen Lange,
``Lengths of Roots of Polynomials in a Hahn Field,''
\emph{Algebra and Logic} \textbf{60} (2021), 95--107.
\href{https://doi.org/10.1007/s10469-021-09632-0}{doi:10.1007/s10469-021-09632-0}.

\end{thebibliography}
\end{document}
